\part{Задание по функциональному анализу (частично).}
\section{Первое задание.}
\subsection{Топологические пространства.}
\begin{task}[T1]
    Рассмотрим семейство подмножество числовой прямой $\tau$, определённое как
    \[
        V \in \tau \Longleftrightarrow V = \emptyset \vee |\R \setminus V| \leq |\N|.
    \]
    \begin{enumerate}
        \item Показать, что $\tau$ -- топология в $\R$.
        \item Найти топологически и секвенциально замкнутые множества в $(\R, \tau)$.
        \item Найти секвенциальное и топологическое замыкания любого более чем счётного множества в $(\R, \tau)$.
        \item Является ли $(\R, \tau)$ хаусдорфовым?
        \item Выполнена ли для $(\R, \tau)$ аксиома счётности?
    \end{enumerate}
\end{task}
\begin{solve}
    \begin{enumerate}
        \item Проверим аксиомы топологического пространства для $\tau$.
        \begin{itemize}
            \item Пустое множество лежит в $\tau$ по определению, $\R \setminus \R = \emptyset$ и не более чем счётно, а следовательно и $\R$ лежит в $\tau$.
            \item Пусть $\{U_\alpha\}_{\alpha \in I} \subset \tau$. Рассмотрим $V \coloneq \bigcup_{\alpha \in I} U_\alpha$. Если все $U_\alpha$ пустые, то $V = \emptyset$ и лежит в $\tau$. Поскольку пустые множества не влияют на объединения, то будем без ограничения общности считать все $U_\alpha$ непустыми. Рассмотрим дополнение $V\colon$
            \[
                \R \setminus V = \R \setminus \bigcup_{\alpha \in I} U_\alpha = \bigcap_{\alpha \in I} \R \setminus U_\alpha.
            \]
            Поскольку все $U_\alpha$ -- непусты, то их дополнения не более чем счётны. В частности, хотя бы для одного $U_\alpha$
            \[
                \bigcap_{\alpha \in I} \R \setminus U_\alpha \subseteq \R\setminus U_{\alpha}.
            \]
            Таким образом $\R \setminus V$ не более чем счётно и лежит в $\tau$ по определению.
            \item Пусть $\{U_i\}_{i = 1}^n \subset \tau$. Рассмотрим $V \coloneq \bigcap_{i = 1}^n U_i$. Если хотя бы один $U_i$ пуст, то и $V = \emptyset$ и лежит в $\tau$ по определению. Пусть все $U_i$ непусты. Рассмотрим дополнение $V\colon$
            \[
                \R \setminus V = \R \setminus \bigcap_{i = 1}^n U_i = \bigcup_{i = 1}^n \R \setminus U_i.
            \]
            Поскольку конечное объединение не более чем счётных множеств само не более чем счётно, то $V \in \tau$ по определению.
        \end{itemize}
        \item Найдём все секвенциально и топологически замкнутые множества в $(\R, \tau)$.
        \begin{itemize}
            \item По определению, замкнутыми множествами являются дополнения открытых. Таким образом замкнутыми являются не более чем счётными множества и $\R$.
            \item Пусть $V \subset \R$ и $\{x_n\} \subset V$ такая, что $x_n \xrightarrow{\tau} x, n \rightarrow \infty$. По определению для любой окрестности $U(x)$ существует $N \in \N$ такой, что $\forall n > N \hookrightarrow x_n \in U(x)$.
            Пусть $K \coloneq \{x_n\} \setminus \{x\}$. Множество $K$ -- не более чем счётно, а следовательно его дополнение открыто. Более того -- его дополнение является окрестностью $x$, а значит начиная с какого-то номера $N$ все элементы последовательности лежат в $\R \setminus K$. Это возможно только если $K$ -- конечно и таким образом $\{x_n\}$ -- стационарная последовательность и, в частности, $x \in V$. Значит в $(\R, \tau)$ всякое подмножество $\R$ секвенциально замкнуто.
        \end{itemize}
        \item Найдём замыкания (топологическое и секвенциальное) для $U \subset \R$ такого, что $|U| > |\N|$.
        \begin{itemize}
            \item Топологическое замыкание $U$ это, по определению, пересечение всех замкнутых, содержащих $U$. В силу того, что все замкнутые множества (кроме $\R$) не более чем счётны, то существует единственное замкнутое множество, содержащее $U$ -- $\R$ и тогда $\overline{U} = \R$.
            \item Как показано в предыдущем пункте, в $(\R, \tau)$ все сходящиеся последовательности стационарны. В частности, из этого следует что всякое множество содержит все свои секвенциальные точки прикосновения и секвенциальные замыканием множества является оно же.
        \end{itemize}
        \item Пусть $(\R, \tau)$ является хаусдорфовым. Тогда для $u, v \in \R\colon u \neq v$ существуют окрестности $U(u)$ и $V(v)$ такие, что $U \cap V = \emptyset$. Поскольку $U$ и $V$ -- непусты, то, из определения топологии, $\R \setminus U$ и $\R \setminus V$ не более чем счётны. Но, тогда, с одной стороны $\R \setminus (U \cap V) = \R \setminus \emptyset = \R$. А с другой -- $\R \setminus (U \cap V) = \R \setminus U \cup \R \setminus V$. Мы приходим к противоречию -- получается, что $\R$ одновременно и не более чем счётно и континуально. Значит $(\R, \tau)$ не является хаусдорфовым пространством.
        \item Для $(\R, \tau)$ не выполнена аксиома счётности: в first-countable пространствах топологическое замыкание совпадает с секвенциальным -- что очевидно не выполняется в $(\R, \tau)$.
    \end{enumerate}
\end{solve}
\begin{task}[T2]
    Пусть $(X, \tau)$ -- топологическое пространство. Рассмотрим следующее семейство подмножеств $X\colon$
    \[
        \mathcal{F} = \{S \subset X \mid \forall \{x_n\} \subset S\colon x_n \xrightarrow{\tau} x, n \rightarrow \infty \hookrightarrow x \in S\}.
    \]
    \begin{enumerate}
        \item Показать, что $\tau_{seq} = \{X \setminus K \mid K \in \mathcal{F}\}$ является топологией.
        \item Сравнить топологии $\tau$ и $\tau_{seq}$.
        \item Показать, что сходимость в $(X, \tau)$ эквивалентна сходимости в $(X, \tau_{seq})$.
    \end{enumerate}
\end{task}
\begin{solve}
    \begin{enumerate}
        \item Покажем что $\tau_{seq}$ является топологией.
        \begin{itemize}
            \item Поскольку всякое замкнутое множество является, в частности, секвенциально замкнутым, то $\emptyset$ и $\R$ лежат в $\tau_{seq}$.
            \item Пусть $\{K_\alpha\}_{\alpha \in I} \subset \mathcal{F}$. Покажем, что $S \coloneq \bigcap_{\alpha \in I} K_\alpha$ лежит в $\mathcal{F}$. Пусть $\{x_n\} \subset S$ такая, что $x_n \xrightarrow{\tau} x, n \rightarrow \infty$. Тогда, в частности, $\forall \alpha \in I \ \{x_n\} \subset K_\alpha$. В силу секвенциальной замкнутости $K_\alpha$ справедливо включение $x \in K_\alpha$. Поскольку оно верно при произвольном $\alpha$, то $x \in S$ и $S$ -- секвенциально замкнуто. Таким образом $S \in \mathcal{F}$.
            \item Пусть $\{K_i\}_{i = 1}^n \subset \mathcal{F}$. Покажем, что $S \coloneq \bigcup_{i = 1}^n K_i$ лежит в $\mathcal{F}$. Пусть $\{x_k\} \subset S$ и $x_k \xrightarrow{\tau} x, k \rightarrow \infty$. Заметим, что существует подпоследовательность $\{x_{k_j}\}$ и номер $s \in \{1, \ldots, n\}$ такой, что $\{x_{k_j}\} \subset K_s$ -- в противном случае она имеет лишь конечное число элементов. Она тоже сходится к $x$. Поскольку $K_s$ -- секвенциально замкнуто, то $x \in K_s$. Значит $x \in S$ и $S$ -- секвенциально замкнуто.
        \end{itemize}
        \item Поскольку каждое замкнутое множество является секвенциально замкнутым, то всякое открытое множество открыто и в секвенциальной топологии, а значит $\tau \subseteq \tau_{seq}$.
        \item Пусть $\{x_n\} \subset X$ и $x \in X$.
        \begin{itemize}
            \item Предположим, что $x_n \xrightarrow{\tau} x, n \rightarrow \infty$.
            То есть
            \[
                \forall U(x) \in \tau \ \exists N \in \N \colon \forall n \geq N \hookrightarrow x_n \in U(x).
            \]
            Пусть, при этом, в $(X, \tau_{seq})$, $x_n  \not\rightarrow x, n \rightarrow \infty$.
            Тогда
            \[
                \exists W(x) \in \tau_{seq} \ \forall N \in \N\hookrightarrow \exists n > N\colon x_n \notin W(x).
            \]
            А значит существует подпоследовательность точек $\{x_{n_j}\} \subset X \setminus W$. В тоже время, она сходится к $x$ в $(X, \tau)\colon$
            \[
                x_{n_j} \xrightarrow{\tau} x, j \rightarrow \infty.
            \]
            По определению $\tau_{seq}$ множество $X \setminus W$ секвенциально замкнуто, а значит $x \in X \setminus W$ -- противоречие: $x \in W$ и, одновременно, $x \in X \setminus W$. Таким образом $x_n \xrightarrow{\tau_{seq}} x, n \rightarrow \infty$.
            \item Предположим, что $x_n \xrightarrow{\tau_{seq}} x, n \rightarrow \infty$. Тогда, по определению сходимости:
            \[\forall U(x) \in \tau_{seq}\  \exists N \in \N \colon \forall n \geq N \hookrightarrow x_n \in U(x).\]
            В тоже время, $\forall U(x) \in \tau \hookrightarrow U(x) \in \tau_{seq}$, а следовательно
            \[
                \forall U(x) \in \tau \hookrightarrow U(x) \in \tau_{seq} \Rightarrow \exists N \in \N \colon \forall n \geq N \hookrightarrow x_n \in U(x).
            \]
            и $x_n \xrightarrow{\tau} x, n \rightarrow \infty$.
        \end{itemize}
    \end{enumerate}
\end{solve}
\begin{task}[T3]
    Пусть $(\R, \tau)$ -- топологическое пространство из задачи T1.
    \begin{enumerate}
        \item Для $\tau$ найти соответствующую топологию $\tau_{seq}$.
        \item Является ли $(\R, \tau_{seq})$ хаусдорфовым?
        \item Выполняется ли для $(\R, \tau_{seq})$ аксиома счётности?
    \end{enumerate}
\end{task}
\begin{solve}
    \begin{enumerate}
        \item Как показано в T1, всякое подмножество $\R$ является секвенциально замкнутым в $(\R, \tau)$, а значит $\tau_{seq} = 2^{\R}$.
        \item Всякая точка в дискретном топологическом пространстве открыта и, следовательно, для $u, v \in \R\colon u \neq v$ мы можем в качестве непересекающихся окрестностей взять просто $\{u\}$ и $\{v\}$.
        \item Да, выполняется. Более того, локальная база окрестностей точки состоит только из одной окрестности -- множества, состоящего из этой точки.
    \end{enumerate}
\end{solve}
\begin{task}[T4]
    В топологическом пространстве $\R^{[0, 1]}$ сравнить секвенциальное и топологическое замыкания множества $C[0, 1]$, состоящего из непрерывных на $[0, 1]$ функций.
\end{task}
\begin{solve}
    Пусть $x \in \R^{[0, 1]}$ и $W(x)$ -- его произвольная окрестность. По определению базы топологии существует такое множество $U(x) \subset W(x)$ такое, что $U(x)$ лежит в базе топологии. Без ограничения общности мы можем считать что $U(x)$ имеет следующий вид:
    \[
        U(x) = \{h \in \R^{[0, 1]} \mid \{t_j\}_{j = 1}^k \subset [0, 1]\colon t_1 < \ldots < t_k \text{ и } \{\epsilon_j\}_{j = 1}^k \subset \R\colon \forall j \hookrightarrow |h(t_j) - x(t_j)| < \epsilon_j\}.
    \]
    Тогда в $U(x)$ лежит $\tilde{h}$ -- кусочно-линейная функция, построеная на разбиении $\{0\} \cup\{t_j\} \cup \{1\}$, и принимающая значения $h(t_j) = x(t_j)$ и произвольные в 0 и 1 (в случае если $t_1 = 0$, то $h(0) = x(0)$; аналогично и с 1). Заметим, что $\tilde{h}$ -- непрерывна, и следовательно $C[0, 1] \cap W(x) \neq \emptyset$. Поскольку $W(x)$ -- произвольная окрестность, то $x$ -- точка прикосновения $C[0, 1]$. Поскольку $x$ выбрана произвольно, то всё $\R^{[0, 1]}$ является точками прикосновения $C[0, 1]$ и, следовательно, топологическим замыканием $C[0, 1]$ является всё пространство. \\
    Известно, что $\R^{[0, 1]}$ со стандартной топологией произведения имеет топологию <<поточечной сходимости>>. В частности, секвенциальное замыкание $C[0, 1]$ состоит из всех функций, которые можно представить в виде поточечного предела непрерывных функций. Заметим, что характеристическая функция множества Витали на отрезке $[0, 1]$ -- $\chi_V$ -- вполне себе функция и лежит в $\R^{[0, 1]} = \overline{C[0,  1]}$. При этом она не является измеримой, поскольку множество Витали -- неизмеримо по Лебегу. С другой стороны, поточечный предел непрерывных функций является измеримой функцией.
    Значит $\chi_V$ не лежит в секвенциальном замыкании $C[0, 1]$ и оно строго меньше топологического.
\end{solve}
\begin{task}[T5]
    Пусть $\tau$ -- топология на $C[0, 1]$, индуцированная из $\R^{[0, 1]}$. Необходимо исследовать отображение $f\colon (C[0, 1], \tau) \to \R$ вида
    \[
        x \mapsto \int_{0}^1 \dfrac{x(t)}{1 + |x(t)|}dt.
    \]
    \begin{enumerate}
        \item ...на топологическую непрерывность.
        \item ...на секвенциальную непрерывность.
    \end{enumerate}
\end{task}
\begin{solve}
    \begin{enumerate}
        \item Покажем топологическую разрывность отображения $f$ в 0. Для этого достаточно показать что существует $\delta > 0$ такая, что во всякой окрестности $0$ в $C[0, 1]$ существует точка $h$ такая, что $f(h) \notin U_\delta(f(0)) = U_\delta(0)$, то есть что $|f(h)| \geq \delta$. \\
        Рассмотрим произвольную окрестность $U(0)$. В ней содержится окрестность из базы $W(0)$, которая, без ограничения общности, имеет вид
        \[
            W(0) = \{h \in C[0, 1] \mid \{t_i\}_{i = 1}^k \colon t_1 < t_2 < \ldots < t_k \wedge |h(t_i)| < \epsilon\}.
        \]
        Зафиксируем $M > 0$ и $\sigma\colon 0 < 2\sigma < l(\{0\} \cup \{t_i\} \cup \{1\})$ и определим функцию $h_{M, \sigma}$ как $M$ вне $\sigma$-окрестностей точек $t_j$, нулём в $\{t_i\}$ и линейно продолжим её до непрерывной. Таких образом определённая функция $h_{M, \sigma}$ лежит в $W(0)$. Оценим $f(h_{M, \sigma})$ следующим образом:
        \[
            f(h_{M, \sigma}) \geq \dfrac{M}{1 + M} \cdot (1 - 2\sigma k).
        \]
        Эта оценка справедлива, поскольку $h_{M, \sigma} \geq 0$ и мера той части интервала, на которой $h_{M, \sigma}$ равен $M$ хотя бы $1 - 2 \sigma k$ (так как у нас $k$ точек $\{t_j\}$ и мера их $\sigma$-окрестностей не более чем $2\sigma k$). \\
        Видно, что
        \[
            \lim\limits_{\overset{M \rightarrow +\infty}{\sigma \rightarrow +0}} f(h_{M, \sigma}) = 1,
        \]
        а следовательно нам достаточно выбрать $\delta < 1$ и для неё мы всегда сможем подобрать $M, \sigma$ такие, чтобы $|f(h_{M, \sigma})|$ была больше $\delta$. \\
        Таким образом $f$ -- топологически разрывна в нуле, а следовательно не является топологически непрерывным отображением.
        \item Секвенциально отображение очевидно непрерывно, поскольку $\forall x(t) \in C[0, 1]\hookrightarrow$
        \[
            \left|\dfrac{x(t)}{1 + |x(t)|}\right| \leq 1.
        \]
        и $1 \in L_1([0, 1], \mathcal{L}^1)$, то по теореме Лебега о мажорируемой сходимости если $x_n \rightarrow x, n \rightarrow \infty$, то
        \[
            \exists \lim\limits_{n \rightarrow \infty} f(x_n) =  \lim\limits_{n \rightarrow \infty} \int_0^1 \dfrac{x_n(t)}{1 +|x_n(t)|}dt = \int_0^1
            \dfrac{x(t)}{1 + |x(t)|}dt = f(x).
        \]
    \end{enumerate}
\end{solve}
\begin{task}[T6]
    Исследовать отображение $f\colon [0, 1]^{[0, 1]} \to [0, 1]$ вида
    \[
        x \mapsto \sum\limits_{k = 1}^{\infty} 2^{-k}x(2^{-k}).
    \]
    \begin{enumerate}
        \item ...на топологическую непрерывность.
        \item ...на секвенциальную непрерывность.
    \end{enumerate}
\end{task}
\begin{solve}
    Покажем что отображение $f$ непрерывно топологически (как следствие -- секвенциально). Рассмотрим следующую последовательность отображений $f_n\colon$
    \[
        f_n = \sum\limits_{k = 1}^n 2^{-k} \pi_{2^{-k}},
    \]
    где под $\pi_k$ понимается натуральная проекция $\pi_k \colon \prod_{i \in [0, 1]} [0, 1] \to [0, 1]_k, \ x \mapsto x(k)$. В силу непрерывности натуральных проекций (тихоновская топология -- слабейшая топология при которой все натуральные проекции непрерывны) всякое отображение $f_n$ непрерывно. В тоже время
    \[
        |f_n| \leq \sum\limits_{k = 1}^n 2^{-k}.
    \]
    И тогда последовательность $f_n$ сходится к $f$ равномерно по признаку Вейерштрасса. Поскольку равномерный предел непрерывных непрерывен, то $f$ -- топологически непрерывное отображение.
\end{solve}
\begin{task}[T7]
    В топологическом пространстве $\R^{(0, 1)}$ рассматривается множество
    \[
        S = \{x \in \R^{(0, 1)} \mid \forall t \in (0, 1) \hookrightarrow t^2 < x(t) < t\}.
    \]
    \begin{enumerate}
        \item Является ли $S$ открытым в $\R^{(0, 1)}$?
        \item Найти топологическое замыкание $S$ в $\R^{(0, 1)}$.
    \end{enumerate}
\end{task}
\begin{solve}
    \begin{enumerate}
        \item Рассмотрим произвольную точку $x \in S$ и её окрестность $U(x)$. Она содержит окрестность из базы $W(x)$, имеющую вид
        \[
            W(x) = \{h \in \R^{(0, 1)} \mid \{t_j\}_{j = 1}^k\colon t_1 < \ldots < t_k \wedge |h(t_i) - x(t_i)| < \epsilon\}.
        \]
        Заметим, что мы можем выбрать точку $t^0 \in (0, 1)$, которая не равна точке из $\{t_{j}\}$ и просто определить $\tilde{x}$ следующим образом:
        \[
            \tilde{x}(t) = \begin{cases}
                               x(t), & t \neq t_0, \\
                               1, & t = t_0.
            \end{cases}
        \]
        Такая $\tilde{x}$ лежит в $W(x)$, но не в $S$, а следовательно никакая точка $S$ не является внутренней (при этом $S$, очевидно, не является $\emptyset\colon \forall t \in (0, 1)$ множество $(t^2, t)$ непусто). Значит $S$ не является открытым множеством в $\R^{(0, 1)}$.
        \item Рассмотрим множество $K$, определяемое как
        \[
            K = \{x \in \R^{(0, 1)} \mid \forall t \in (0, 1) \hookrightarrow t^2 \leq x(t) \leq t\}.
        \]
        Покажем, что всякая точка $K$ является предельной для $S$. Пусть $x \in K$. Рассмотрим произвольную окрестность $U(x)$. Она содержит в себе окрестность из базы $W(x)$, которая имеет вид
        \[
            W(x) = \{h \in \R^{(0, 1)} \mid \{t_j\}_{j = 1}^k\colon t_1 < \ldots < t_k \wedge |h(t_i) - x(t_i)| < \epsilon\}.
        \]
        Пусть $A = \{t \in (0, 1) \mid x(t) = t\}$ и $B = \{t \in (0, 1) \mid x(t) = t^2\}$. Определим функцию $\tilde{x}$ следующим образом:
        \[
            \tilde{x}(t) = \begin{cases}
                               t - \frac{\epsilon}{2}, & t \in A \wedge \exists j\colon t = t_j, \\
                               t^2 + \frac{\epsilon}{2}, & t \in B \wedge \exists j\colon t = t_j, \\
                               x(t), & \text{иначе}.
            \end{cases}
        \]
        (Без ограничения общности будем считать $\epsilon$ достаточно маленьким, чтобы при всяких $t_j \in (0, 1)$ было верно, что $t_j - t_j^2 > \epsilon$). Определённая таким образом функция $\tilde{x}$ лежит и в $S$, и в $W(x)$, то есть $U(x) \cap S \neq \emptyset$. Значит, всякая точка $K$ является предельной для $S$. \\
        Покажем теперь, что $\R^{(0, 1)} \setminus K$ -- открыто. Пусть $x \in K^C$, то есть
        \[
            \exists t_0 \in (0, 1)\colon x(t_0) < t_0^2 \vee x(t_0) > t_0.
        \]
        Определим окрестность $U(x)$ следующим образом:
        \[
            U(x) = \{h \in \R^{(0, 1)} \mid |h(t_0) - x(t_0)| < \epsilon\},
        \]
        где $\epsilon = (t_0^2 - x(t_0)) / 2$, если $x(t_0) < t_0^2$ и $\epsilon = x(t_0) - t_0$ иначе. Определённое таким образом $U(x)$ открыто по определению тихоновской топологии и имеет пустое пересечение с $K$, а следовательно $x$ -- внутренняя точка $K^C$. В силу произвольности выбора $x$, $K^C$ -- открыто и $K$ -- замкнуто. \\
        Поскольку $S \subset K$ и $K$ -- замкнуто, то $\overline{S} \subset K$. С другой стороны, всякая точка $K$ -- предельна для $S$, а значит $\overline{S} \supset K$. Таким образом $K = \overline{S}$.
    \end{enumerate}
\end{solve}
\begin{task}[T8]
    В топологическом пространстве $\R^{\N}$ рассматривается множество
    \[
        S = \{x \in \R^{\N} \mid \forall n \in \N \hookrightarrow n < x(n) < 2^n\}.
    \]
    \begin{enumerate}
        \item Является ли $S$ открытым в $\R^{\N}$?
        \item Найти топологическое замыкание $S$ в $\R^{\N}$.
    \end{enumerate}
\end{task}
\begin{solve}
    \begin{enumerate}
        \item Пусть $x \in S$. Рассмотрим окрестность $U(x)$.

        В ней лежит некоторая окрестность из базы $W(x)$, которая без ограничения общности имеет вид
        \[
            W(x) = \{h \in \R^{\N} \mid \{t_j\}_{j = 1}^k \subset \N \colon t_1 < \ldots < t_k \wedge |h(t_j) - x(t_j)| < \epsilon\}.
        \]
        Существует $N \in \N$ такой, что $\forall j \in \{1, \ldots, k\} \hookrightarrow N > t_j$. Определим $\tilde{x}$ следующим образом:
        \[
            \tilde{x}(m) = \begin{cases}
                               2^m, & m = N \\
                               x(m), & m \neq N.
            \end{cases}
        \]
        Такая функция $\tilde{x}$ лежит в $W(x)$, но не лежит в $S$, а следовательно $U(x) \not\subset S$ -- и никакая точка $S$ не является внутренней, а значит $S$ не является открытым.
        \item Заметим, что $\R^{\N}$ -- счётное произведение пространств, удовлетворяющих аксиоме счётности, а следовательно и $\R^{\N}$ удовлетворяет аксиоме счётности.

        В таком случае топологическое замыкание эквивалентно секвенциальному. Покажем, что
        \[
            K = \{x \in \R^{\N} \mid \forall n \in \N \hookrightarrow n \leq x(n) \leq 2^n\}.
        \]
        является топологическим замыканием $S$. \\
        Пусть $x$ -- предельная точка $S$. Тогда существует последовательность $\{x_n\} \subset S$ такая, что $x_n \rightarrow x, n \rightarrow \infty$. Но
        \[
            m < x_n(m) < 2^m.
        \]
        Перейдём к пределу по $n$ и получим:
        \[
            m \leq x(m) \leq 2^m.
        \]
        И таким образом $x \in K$, а значит $\overline{S} \subset K$. \\
        Пусть $x$ -- точка из $K$. Определим последовательность точек $S$, сходящуюся к $x$. Пусть $A = \{m \in \N \mid x(m) = m\}$ и $B = \{m \in \N \mid x(m) = 2^m\}$.
        Определим $x_n$ следующим образом:
        \[
            x_n(m) = \begin{cases}
                         x(m) + \frac{1}{2n}, & m \in A \\
                         x(m) - \frac{1}{2n}, & m \in B \\
                         x(m), & m \notin A \cup B.
            \end{cases}
        \]
        При всех $n \in \N$ отображение $x_n$ лежит в $S$, а при $n \rightarrow \infty$ сходится к $x \in K$ -- значит, всякая точка $K$ -- предельна для $S$ и $\overline{S} \supset K$. \\
        Таким образом $\overline{S} = K$.
    \end{enumerate}
\end{solve}
\begin{task}[T9]
    Привести пример множества из $\R^{[0, 1]}$, секвенциальное замыкание которого не совпадает с его топологическим замыканием в пространстве $\R^{[0, 1]}$.
\end{task}
\begin{solve}
    $C[0, 1]$. Как показано в T4, его топологическим замыканием является $\R^{[0, 1]}$, а секвенциальное строго меньше.
\end{solve}
\begin{task}[T10]
    Привести пример множества из $\R^{[0, 1]}$, секвенциальное замыкание которого в пространстве $\R^{[0, 1]}$ не является секвенциально замкнутым.
\end{task}
\begin{solve}
    И снова $C[0, 1]$. Для доказательства этого нам понадобятся три факта:
    \begin{enumerate}
        \item Поточечный предел непрерывных функций на $[0, 1]$ имеет всюду плотное множество в $[0, 1]$ точек непрерывности.
        \item Индикатор точки является пределом непрерывных функций.
        \item Функция Дирихле $\chi_{\Q \cap [0, 1]}$ является пределом суммы индикаторов точек.
    \end{enumerate}
    По модулю этих трёх фактов мы можем сказать что $\chi_{\Q \cap [0, 1]}$ не принадлежит секвенциальному замыканию $C[0, 1]$ -- она всюду разрывна. С другой стороны, она принадлежит второму секвенциальному замыканию $C[0, 1]$, поскольку $\chi_{\Q \cap [0, 1]} = \sum_{q \in \Q \cap [0, 1]} \chi_{\{q\}}$, а каждый из индикаторов точки является пределом непрерывных функций. Таким образом секвенциальное замыкание $C[0, 1]$ не является секвенциально замкнутым: $\chi_{\Q \cap [0, 1]}$ является его секвенциальной предельной точкой, но не лежит в нём. \\
    Теперь покажем справедливость этих трёх утверждений.
    \begin{enumerate}
        \item Рассмотрим $\{x_m\} \subset C[0, 1]$ такую, что $x_m \rightarrow x \in \R^{[0, 1]}$.

        Определим несколько последовательностей множеств:
        \[P_{m, n} = \{t \in [0, 1] \mid |x_m(t) - x(t)| \leq \frac{1}{n}\}, \quad Q_n = \bigcup_{m = 1}^{\infty} \Int P_{m, n}, \quad R = \bigcap_{n = 1}^{\infty} Q_n.\]
        Пусть $t \in R$. Тогда, в частности, $\forall n \in \N\colon t \in Q_n$. Поскольку \[t \in Q_n \Longleftrightarrow \exists m\in \N \exists \delta > 0\colon \forall s \in U_\delta(t) \hookrightarrow |x_m(s) - x(s)| \leq \frac{1}{n}.\]
        То есть
        \[
            t \in R \Longleftrightarrow \forall n \in \N \exists m \in \N \exists \delta > 0\colon \forall s \in U_\delta(t) \hookrightarrow |x_m(s) - x(s)| \leq \frac{1}{n.}
        \]
        Тогда, в частности, если $t \in R$, то
        \[\forall \epsilon > 0 \ \exists N = \lceil \frac{1}{\epsilon} \rceil \  \exists m = m(N) \ \exists \delta = \delta(N) > 0\colon  \forall s \in U_{\delta}(t) \hookrightarrow |x_m(s) - x(s)| \leq \frac{1}{N} \leq \epsilon.\]
        Поскольку $x_m$ -- непрерывна, то
        \[
            \forall \epsilon > 0 \ \exists \delta_0 > 0 \ \forall s \in U_{\delta_0}(t) \hookrightarrow |x_m(s) - x_m(t)| \leq \epsilon.
        \]
        Тогда в $\min\{\delta, \delta_0\}$-окрестности справедливы обе оценки и тогда, в силу неравенства треугольника, для $s \in U_{\min\{\delta_0, \delta\}}(t)\colon$
        \[
            |x(s) - x(t)| \leq |x(s) - x_m(s)| + |x_m(s) - x_m(t)| + |x_m(t) - x(t)| \leq 3\epsilon.
        \]
        Тогда $t$ -- точка непрерывности $x$.

        Пусть $t$ -- произвольная точка непрерывности $x$. Тогда
        \[
            \forall \epsilon > 0 \ \exists \delta > 0 \ \forall s \in U_\delta(t) \hookrightarrow |x(t) - x(s)| \leq \epsilon. \tag{1}
        \]
        С другой стороны, поскольку $x_m(t) \rightarrow x(t), m \rightarrow \infty$, то
        \[
            \forall \epsilon > 0 \ \exists N \in \N \  \forall m \geq N \hookrightarrow |x_m(t) - x(t)| \leq \epsilon. \tag{2}
        \]
        Также, в силу непрерывности $x_m\colon$
        \[
            \forall \epsilon > 0 \ \exists \delta_0 > 0 \  \forall s \in U_{\delta_0}(t) \hookrightarrow |x_m(s) - x_m(t)| \leq \epsilon. \tag{3}
        \]
        Тогда $\forall n \in \N \ \exists \epsilon > 0 \colon 3\epsilon \leq \frac{1}{n}$.  Из условия $(2)$ существует $m = N(\epsilon)$, а из условий $(1)$ и $(3)$ мы можем выбрать $\delta_1 = \min\{\delta, \delta_0\}$ для которых верно
        \[
            \forall s \in U_{\delta_1}(t) \hookrightarrow |x_m(s) - x(s)| \leq |x_m(s) - x_m(t)| + |x_m(t) - x(t)| + |x(t) - x(s)| \leq 3\epsilon \leq \frac{1}{n}.
        \]
        Что и показывает принадлежность точки $t$ множеству $R$. Таким образом $R$ -- $G_\delta$-множество которое состоит в точности из точек непрерывности функции $x$.

        Покажем плотность $Q_n$ в $[0, 1]$ при каждом $n \in \N$. Пусть $Q_n$ -- не плотно. Тогда существует такой интервал $(a, b) \subset [0, 1]$, что $(a, b) \cap Q_n = \emptyset \Longleftrightarrow \forall m \in \N\colon \Int P_{m, n} \cap (a, b) = \emptyset$.

        Рассмотрим следующие множества:
        \[
            A_{m, n} = \bigcap_{k,j\geq m}\{t \in [0, 1] \mid |x_j(t) - x_k(t)| \leq \frac{1}{n}\}  .
        \]
        Заметим, что каждое $A_{m, n}$ -- замкнуто как пересечение замкнутых (которые замкнуты в силу непрерывности $x_m$). Так как $x_m \rightarrow x, m \rightarrow \infty$, то последовательность $\{x_m\}$ -- фундаментальна и $\forall t \in [0, 1] \  \exists m \in \N\colon \forall k,j \geq m \hookrightarrow |x_j(t) - x_k(t)| \leq \frac{1}{n}$. Тогда
        \[
            [0, 1] = \bigcup_{m = 1}^{\infty} A_{m, n}.
        \]
        Также, если $t \in A_{m, n}$, то в силу поточечной сходимости $\frac{1}{n} \geq |x_j(t) - x_k(t)| \rightarrow |x_j(t) - x(t)|$ при $j \geq m$, а значит $t \in P_{m, n}$. Тогда  $\Int A_{m, n} \subset \Int P_{m, n}$.

        Значит $\forall m \in \N\colon \Int A_{m, n} \cap (a, b) = \emptyset$. В частности, если $[c, d] \subset (a, b)$ и $c \neq d$, то
        \[
            [c, d] = \bigcup_{m = 1}^{\infty} A_{m, n} \cap [c, d].
        \]
        Но $A_{m, n}$ -- замкнутые множества, не имеющие внутренности в $[c, d]$, а следовательно мы получили представление полного метрического пространства в виде счётного объединения нигде не плотных множеств -- противоречие. Следовательно, каждое $Q_n$ -- всюду плотно, а следовательно и $R$ всюду плотно, что и требовалось доказать.

        \item Пусть $q \in [0, 1]$. Определим $x_n$ как
        \[
            x_n(t) =
            \begin{cases}
                0,  & t \leq q - \frac{1}{2n}, \\
                1 - 2n(q - t), & q - \frac{1}{2n} \leq t \leq q, \\
                1 + 2n(q - t), & q \leq t \leq q + \frac{1}{2n}, \\
                0, & q + \frac{1}{2n} \leq t.
            \end{cases}
        \]
        Очевидно, что $x_n \rightarrow \chi_{q}, n \rightarrow \infty$. Таким образом $\chi_{q}$ лежит в секвенциальном замыкании $C[0, 1]$ при всяком $q \in [0, 1]$ (в частности, при $q \in \Q \cap [0, 1]$).
        \item Рассмотрим некоторую нумерацию $\{q_k\}$ рациональных чисел на отрезке $[0, 1]$. Определим $x_n = \sum\limits_{k = 1}^n \chi_{q_k}$. Тогда $x_n \rightarrow \chi_{\Q \cap [0, 1]}, n \rightarrow \infty$.

    \end{enumerate}
\end{solve}
\begin{task}[T12]
    Рассмотрим множество $\mathfrak{M}[0, 1]$ из задачи T11. Будем говорить, что множество $S \subset \mathfrak{M}[0, 1]$ секвенциально замкнуто (относительно сходимости почти всюду), если для любой последовательности $\{x_n\} \subset S$ такой, что $x_n \xrightarrow{a.s.} x \in \mathfrak{M}[0, 1]$ выполняется $x \in S$.
    \begin{enumerate}
        \item Доказать, что семейство
        \[
            \tau_\mu = \{\mathfrak{M}[0, 1] \setminus S \mid S \text{ -- секвенциально замкнуто относительно сходимости п.в. }\}.
        \]
        является топологией в $\mathfrak{M}[0, 1]$.
        \item Доказать, что для последовательности $\{x_n\} \subset \mathfrak{M}[0, 1]$ и функции $x \in \mathfrak{M}[0, 1]$ выполнено
        \[
            x_n \xrightarrow{\tau_\mu} x \Longleftrightarrow x_n \xrightarrow{\mu} x,
        \]
        где под <<$\xrightarrow{\mu}$>> понимается сходимость по мере:
        \[
            \forall \epsilon > 0 \ \lim\limits_{n \rightarrow \infty} \mu\left\{t \in [0, 1] \mid |x_n(t) - x(t)| > \epsilon\right\} = 0.
        \]
    \end{enumerate}
\end{task}
\begin{solve}
    \begin{enumerate}
        \item Доказательство того, что $\tau_\mu$ -- топология полностью аналогично доказательству того, что $\tau_{seq}$ является топологией.
        \item Пусть $\{x_n\} \subset \mathfrak{M}[0, 1]$ и $x \in \mathfrak{M}[0, 1]$.
        \begin{itemize}
            \item Пусть $x_n \xrightarrow{\mu} x, n \rightarrow \infty$. Покажем, что
            \[
                \forall V\colon V^C \in \tau_\mu \wedge x \notin V \ \exists N \in \N\colon \forall n \geq N \hookrightarrow x_n \notin V.
            \]
            Пусть существует такое секвенциально замкнутое относительно сходимости почти всюду множество $V$, не содержащее $x$, что $\forall N \in \N \exists n > N\colon x_n \in V$. Тогда у нас есть подпоследовательность $\{x_n\}$, лежащая в $V$ и сходящаяся по мере к $x$. По теореме Рисса из неё можно выделить подпоследовательность, сходящуюся почти всюду и тогда, в силу секвенциальной замкнутости, $x \in V$ -- противоречие. Следовательно, такого множества не существует и $x_n \xrightarrow{\tau_\mu} x, n \rightarrow \infty$.
            \item Пусть $x_n \xrightarrow{\tau_\mu} x, n \rightarrow \infty$, но не сходится по мере, то есть
            \[
                \exists \epsilon > 0 \  \exists \{x_{n_k}\} \  \exists \delta > 0 \colon \mu\{t \in [0, 1] \mid |x_{n_k}(t) - x(t)| \geq \epsilon\} \geq \delta.
            \]
            Обозначим как $A_{k} = \{t \in [0, 1] \mid |x_{n_k} - x(t)| \geq \epsilon\}$ и рассмотрим $B = \limsup\limits_{k \rightarrow \infty} A_{k}$. Заметим, что $\mu(B) \geq \delta > 0$, поскольку $\mu(A_k) \geq \delta$ (в силу непрерывности меры). Тогда никакая подпоследовательность $\{x_{n_k}\}$ не может сходиться почти всюду к $x$, поскольку на множестве положительной меры $B$ отсутствует сходимость. Рассмотрим множество $S$ -- секвенциальное замыкание $\{x_{n_k} \mid k \in \N\}$. Заметим, что $S$ -- секвенциально замкнуто (в силу диагонализируемости сходимости почти всюду на пространстве $\sigma$-конечной меры), содержит $\{x_{n_k}\}$ (можем взять постоянные последовательности) и не содержит $x$ (как показано выше). Но, тогда, $S^C$ -- открытое в $\tau_\mu$ множество, являющееся окрестностью $x$ и не содержащее бесконечное число элементов последовательности, что вступает в противоречие со сходимостью по $\tau\mu$ к $x$. Таким образом $x_n \xrightarrow{\mu} x, n \rightarrow \infty$.
        \end{itemize}
    \end{enumerate}
\end{solve}
\begin{task}[T11]
    Рассмотрим множество функций
    \[
        \mathfrak{M}[0, 1] = \{x\colon[0, 1] \to \R \mid x \text{ -- измерима по Лебегу }\}.
    \]
    Функции из $\mathfrak{M}[0, 1]$, совпадающие почти всюду, отождествим. Рассмотрим в $\mathfrak{M}[0, 1]$ сходимость почти всюду последовательности $\{x_n\} \subset \mathfrak{M}[0, 1]$ к $x \in \mathfrak{M}[0, 1]\colon$
    \[
        x_n \xrightarrow{a.s.} x \Longleftrightarrow \mu\left\{t \in [0, 1] \mid \exists \lim\limits_{n \rightarrow \infty} x_n(t) = x(t)\right\} = 1.
    \]
    Существует ли топология $\tau$ в $\mathfrak{M}[0, 1]$ такая, что
    \[
        x_n \xrightarrow{a.s.} x \Longleftrightarrow x_n \xrightarrow{\tau} x?
    \]
\end{task}
\begin{solve}
    Предположим, что существует такая $\tau$, что $x_n \xrightarrow{a.s.} x, n \rightarrow \infty \Longleftrightarrow x_n \xrightarrow{\tau} x, n \rightarrow \infty$. Заметим, что в таком случае $\tau_\mu$ -- это в точности $\tau_{seq}$ для $\tau$. Поскольку сходимость по $\tau$ и $\tau_{seq}$ эквивалентна, то
    \[
        x_n \xrightarrow{a.s.} x \Longleftrightarrow x_{n} \xrightarrow{\tau} x \Longleftrightarrow x_n \xrightarrow{\tau_\mu} x \Longleftrightarrow x_n \xrightarrow{\mu} x.
    \]
    Что очевидно неверно: возьмём <<бегающий уменьшающийся индикатор>> (то есть последовательность функций вида $f_n = \chi_{[k/2^m, (k + 1)/2^m]}$)). Такая последовательность сходится по мере к нулю, но всюду поточечно расходится.
\end{solve}
\subsection{Метрические пространства.}
\begin{task}[\S 1.3]
    Является ли открытым в пространстве $C[a, b]$ множество
    \[
        S = \{f \in C[a, b] \mid \forall x \in [a, b] \hookrightarrow 0 < f(x) < 1\}.
    \]
\end{task}
\begin{solve}
    Рассмотрим $f \in S$. Образ связного замкнутого компакта $[a, b]$ под действием непрерывного отображения $f$ -- это связный замкнутый компакт $[m, M] \subset (0, 1)$. Пусть $\epsilon = \min\left\{\frac{m}{2}, \frac{1 - M}{2}\right\}$. Покажем, что $B_\epsilon(f) \subset S$. Пусть $g \in B_\epsilon(f)$. Тогда, так как
    \[
        g = f + (g - f),
    \]
    то
    \[
        \min g \geq \min f + \min (g - f) = m + \min(g - f) \geq m - \max |f - g| \geq m - \epsilon \geq \frac{m}{2} > 0.
    \]
    и
    \[
        \max g \leq \max f + \max (g - f) = M + \max (g - f) \leq M + \epsilon \leq \frac{1 + M}{2} < 1.
    \]
    А значит $g \in S$ и точка $f$ -- внутренняя. Значит $S$ -- открытое, так как всякая его точка является внутренней.
\end{solve}
\begin{task}[\S 1.4]
    Является ли открытым в пространстве $l_\infty$ множество
    \[
        S = \{x \in l_\infty \mid \forall k \in \N \hookrightarrow 0 < x(k) < 1\}.
    \]
\end{task}
\begin{solve}
    Рассмотрим последовательность $x \in S$, которая определяется как $x(n) = \frac{1}{n + 1}$. Она является существенно ограниченной: $\|x\|_\infty = \sup\limits_{n \in \N} |x(n)| = 0.5$. Определённая таким образом последовательность лежит в $S$, поскольку $0 < \frac{1}{n + 1} < 1$ при всех $n \in \N$. С другой стороны, для всякой $\epsilon$-окрестности существует последовательность $x_\epsilon$, которая определяется как
    \[
        x_\epsilon(n) = \begin{cases}
                            x(n), & \frac{1}{n + 1} \geq \epsilon/2, \\
                            0, & \frac{1}{n + 1} < \epsilon/2.
        \end{cases}
    \]
    Такая последовательность очевидно лежит в $l_\infty$ и лежит в $B_\epsilon(x)$, поскольку она начинает отличаться только при $\frac{1}{n + 1} < \epsilon/2$, то есть $\sup |x - x_\epsilon| \leq \epsilon/2 < \epsilon$. Значит $x$ -- не внутренняя точка $S$ и, следовательно, $S$ не является открытым.
\end{solve}
\begin{task}[\S 1.5]
    Пусть $A$  -- подмножество метрического пространства $(X, \rho)$. Доказать, что функция $f\colon X \to \R, x \mapsto \rho(x, A) = \inf\limits_{y \in A} \rho(x, y)$ непрерывна.
\end{task}
\begin{solve}
    Пусть $x_n \rightarrow x, n \rightarrow \infty$. Тогда
    \[
        \rho(x, y) \leq \rho(x, x_n) + \rho(x_n, y).
    \]
    Возьмём $\inf$ по $y \in A\colon$
    \[
        f(x) \leq \rho(x, x_n) + f(x_n).
    \]
    А теперь рассмотрим нижний предел по $n \in \N\colon$
    \[
        f(x)  \leq \liminf f(x_n).
    \]
    С другой стороны
    \[
        \rho(x_n, y) \leq \rho(x, x_n) + \rho(x, y).
    \]
    Перейдёс к $\inf$ по $y \in A\colon$
    \[
        f(x_n) \leq \rho(x, x_n) + f(x).
    \]
    И возьмём верхний предел по $n \in \N\colon$
    \[
        \limsup f(x_n) \leq f(x).
    \]
    Тогда $f(x) = \limsup f(x_n) = \liminf f(x_n)$, а значит существует $\lim f(x_n) = f(x)$ и функция расстояния непрерывна.
\end{solve}
\begin{task}[\S 1.7]
    Пусть $A, B$  -- замкнутые, непересекающиеся подмножества метрического пространства $(X, \rho)$. Доказать, что на $X$ существует непрерывная функция $f$ такая, что $f|_{A} \equiv 0, f|_{B} \equiv 1$.
\end{task}
\begin{solve}
    Известно, что функции $\rho_A = \rho(\cdot, A)$ и $\rho_B = \rho(\cdot, B)$ -- непрерывны. Определим $h$ следующим образом:
    \[
        h = \frac{\rho_A}{\rho_A + \rho_B}.
    \]
    Заметим, что $\rho_A + \rho_B > 0$. Пусть $\rho_A(x) + \rho_B(x) = 0$. Тогда существует последовательности $\{z_n\}$ и $\{y_n\}$ такая, что они сходятся к $x$ и лежат в $A$ и $B$ соответственно (в силу определения инфинума). В силу замкнутости $A$ и $B$ точка $x$ лежит и в $A$, и в $B$ -- что вступает в противоречие с тем, что $A$ и $B$ -- непересекающиеся.
    Тогда $h$ -- корректно определённая непрерывная функция, которая и является искомой.
\end{solve}
\begin{task}[\S 1.13]
    Доказать, что пространство основных функций $D(\R^1)$ неметризуемо.
\end{task}
\begin{solve}
    Пусть $w$ -- соболевская функция:
    \[
        w(t) = \begin{cases}
                   e^{\frac{1}{1 - \|x\|^2}}, & \|x\| < 1, \\
                   0, & \|x\| \geq 1.
        \end{cases}
    \]
    Рассмотрим следующее двух-параметрическое семейство функций $\{\psi_{m, n}\}_{m, n \in \N}\colon$
    \[
        \psi_{m, n}(t) = \frac{1}{m}w(t) + \frac{1}{n}w(t - m).
    \]
    Заметим, $\{\psi_{m, n}\} \subset D(\R^1)$.

    При фиксированном $m$ и $n \rightarrow \infty$ последовательность $\psi_{m, n}$ сходится в $D$ к $\frac{1}{m}w(t)$: носители лежат в отрезке $[-1, m + 1]$ а производные равномерно ($k$-ая производная $w$ ограничена некоторой $\frac{C_k}{n}$, где $C_k = \max w^{(k)}$) сходятся к 0 на $[m - 1, m + 1]$, а в остальных точках функция $\{\psi_{m, n}\}_{n = 1}^{\infty}$ совпадают. \\
    Взяв второй предел мы сможем получить $0 \in D(\R^1)\colon \frac{1}{m}w \xrightarrow{D} 0, m \rightarrow \infty$. \\
    Покажем, что 0 не лежит в первом секвенциальном замыкании семейства $\{\psi_{m, n}\}_{m, n \in \N}$ (поскольку во втором он, как показано, лежит). Пусть $\{\phi_k\}_{k = 1}^{\infty} \subset \{\psi_{m, n}\}_{m, n \in \N}$ такое, что $\phi_{k} \xrightarrow{D} 0, k \rightarrow \infty$. Поскольку носители элементов последовательности лежат в некотором компакте $K$, то (поскольку всякое $\phi_k = \psi_{m(k), n(k)}$) $\sup\limits_{k}m(k) < \infty$. Но, тогда, $\phi_k(0) \geq \frac{1}{\sup\limits_k m(k)}w(0) > 0$ и последовательность $\phi_k$ не может сходиться к нулевой функции. \\
    Заметим, что в пространстве, удовлетворяющем первой аксиоме счётности (а всякое метрическое пространство ей удовлетворяет) секвенциальное замыкание совпадает с топологическим и, как следствие, секвенциальное замыкание само по себе секвенциально замкнуто. Таким образом сходимость в $D$ не может быть метризуема.
\end{solve}
\begin{task}[\S 2.2]
    Доказать, что пространства $l_p \  (1 \leq p < \infty)$ -- полные сепарабельные метрические пространства, а $l_\infty$ -- полно, но не сепарабельно.
\end{task}
\begin{solve}
    Пространства последовательностей $l_p$ допускают характеризацию как $L^p(\N, m)$ где $m$ -- считающая мера. Поэтому для доказательства их полноты достаточно показать полноту пространств $L^p(X, \mu)$ (пусть $\mu$ -- полная мера). В силу критерия полноты нормированного пространства нам достаточно доказать сходимость каждого абсолютно сходящегося ряда. \\
    Пусть $\{f_n\} \subset L^p(X, \mu)$ и ряд $\sum\limits_{n = 1}^{\infty} f_n$ сходится абсолютно, то есть $\sum\limits_{n = 1}^{\infty} \|f_n\|_p < +\infty$. Положим $F_k = \left(\sum\limits_{n = 1}^k |f_n|\right)^p$. Последовательность $\{F_k\}$ -- неубывающая и состоит из неотрицательных функций. В силу неравенства Минковского:
    \[
        \left(\int_X F_nd\mu\right)^{1/p} \leq \sum\limits_{k = 1}^n \|f_k\|_p \leq \sum\limits_{k = 1}^{\infty}\|f_k\|_p < +\infty.
    \]
    Таким образом $\forall n \in \N$ функция $F_n$ -- интегрируема.
    По теореме о монотонной сходимости:
    \[
        \exists \lim\limits_{n \rightarrow \infty} \left(\int_X F_n d\mu\right)^{1/p} = \left(\int_X \lim\limits_{n \rightarrow \infty} F_n d\mu\right)^{1/p},
    \]
    А значит
    \[
        \left(\int_X \left(\sum\limits_{n = 1}^{\infty} |f_n|\right)^pd\mu\right)^{1/p} \leq \sum\limits_{n = 1}^{\infty} \|f_n\|_p < +\infty.
    \]
    Тогда и функция $\sum\limits_{n = 1}^{\infty} |f_n|$ -- интегрируема и, следовательно, почти всюду конечна. Таким образом для почти всех $x \in X$ верно $\sum\limits_{n = 1}^{\infty} |f_n(x)| < +\infty$. Значит функция $F = \sum\limits_{n = 1}^{\infty} f_n$ существует почти всюду. Доопределим её 0 в точках, в которых суммы ряда не существует. \\
    Покажем что $\left\|F - \sum\limits_{k = 1}^n f_k\right\|_p \rightarrow 0, n \rightarrow \infty$. Поскольку
    \[
        \left(\int_X \left|\sum\limits_{k = n}^m f_k\right|^p\right)^{1/p} \leq \sum\limits_{k = n}^{m} \|f_k\|_p \leq \sum\limits_{k = n}^{\infty}  \|f_k\|_p.
    \]
    Применение леммы Фату завершает доказательство полноты:
    \[
        \left\|F - \sum\limits_{k = 1}^n f_k\right\|_p \leq \liminf_{m \rightarrow \infty} \left\|\sum\limits_{k = n}^{m} f_k\right\|_p \leq \sum\limits_{k = n}^{\infty} \|f_k\|_p \rightarrow 0, n \rightarrow \infty.
    \]


    Покажем сепарабельность пространств $l_p$ при $p \in [1, \infty)$.
    Заметим, что множество финитных последовательностей $c_0$ плотно в $l_p$. Пусть
    $x \in l_p$ и $x^N$ определяется как
    \[
        x^N(n) = \begin{cases}
                     x(n), & n \leq N \\
                     0, & n > N.
        \end{cases}
    \]
    Всякая $x^N \in c_0$. С другой стороны
    \[
        \|x^N - x\|_p^p = \sum\limits_{k = N + 1}^{\infty} x(k)^p \rightarrow 0, n \rightarrow \infty.
    \]
    С другой стороны, в силу плотности $\Q$ в $\R$ множество финитных последовательностей над $\Q$ -- $c_0^{\Q}$ -- плотно в $c_0$, а значит и в $l_p$. Поскольку $c_0^{\Q}$ -- счётно, то $l_p$ -- сепарабельно.

    Покажем несепарабельность пространства $l_\infty$. Для этого достаточно построить несчётное семейство элементов $l_\infty$ таких, что расстояние между двумя различными элементами семейства будет $\geq 1$. Пусть $\emptyset \neq B \subseteq \N$. Пусть $x_B = \chi_B$. Тогда при $A \neq B$ будет существовать точка $x \in A \triangle B$, а следовательно $\|\chi_A - \chi_B\|_\infty \geq |\chi_A(x) - \chi_B(x)| \geq 1$. Поскольку $|2^{\N}| = |\R|$, то это множество континуально и $l_\infty$ не является сепарабельным.
\end{solve}
\begin{task}[\S 2.3]
    Доказать, что если в пространстве $C[a, b]$ рассмотреть метрику $\rho_1(f, g) = \int_a^b |f - g|dx$, то в ней оно будет неполно.
\end{task}
\begin{solve}
    Рассмотрим последовательность функций $\{f_n\}_{n= 1}^\infty$, определённую следующим образом:
    \[
        f_n(t) =
        \begin{cases}
            0, & a \leq t \leq \frac{a + b}{2} - \frac{1}{2n}, \\
            2n(t - \frac{a + b}{2} + \frac{1}{2n}), & \frac{a + b}{2} - \frac{1}{2n} \leq t \leq \frac{a + b}{2}, \\
            1, & \frac{a + b}{2} \leq t \leq b.
        \end{cases}
    \]
    Она фундаментальна по $L_1$-метрике:
    \[
        \int_a^b |f_n(t) - f_m(t)|dt \leq \left|\frac{1}{4m} - \frac{1}{4n}\right| \leq \dfrac{1}{2m} \xrightarrow[m \rightarrow \infty
        ]{} 0.
    \]
    При этом, при естественном вложении $(C[a, b], \|\cdot\|_1)$ в $L_1([a, b])$ в нём последовательность сходится к $f\colon [a, b] \to \R\colon$
    \[
        f(t) = \begin{cases}
                   1, & t \geq \frac{a + b}{2}, \\
                   0, & t < \frac{a + b}{2}.
        \end{cases}
    \]
    Которая не является непрерывной с точностью до меры нуль, а следовательно пространство $C[a, b]$ с $L_1$-метрикой неполно.
\end{solve}
\begin{task}[\S 2.6]
    Найти пополнение метрического пространства, состоящего из непрерывных финитных на $\R$ функций (то есть $C_0(\R)$) с метрикой
    \[
        \rho(x, y) = \max\limits_{t \in \R} |x(t) - y(t)|.
    \]
\end{task}
\begin{solve}
    Пусть $C_{\pm\infty}(\R)$ -- пространство непрерывных функций с пределом 0 в бесконечности и $\sup$-метрикой.
    \begin{enumerate}
        \item Покажем, что $C_{\pm\infty}(\R)$ -- полно по $\sup$-метрике. Пусть $\{f_n\}$ -- фундаментальная последовательность. Заметим, что она, в частности, фундаментальна в каждой точке и, следовательно, в силу полноты $\R$, существует некоторая функция $f$ такая, что $f_n \rightarrow f, n \rightarrow$ поточечно. В силу фундаментальности по $\sup$-метрике сходимость $f_n$ к $f$ на самом деле равномерная и тогда $f$ -- непрерывна (или же, в силу того, что непрерывность -- локальное свойство, и от сужения области определения на некоторый компакт фундаментальность не теряется, то последовательность сходится равномерно к функции на всяком компакте -- следовательно, предел является непрерывной функции). Покажем, что $f$ имеет нулевой предел на бесконечности. Пусть это не так. Тогда $\exists \epsilon > 0 \colon \forall t > 0 \exists r > t \colon f(r) > \epsilon$. Поскольку $f_n \rightrightarrows f, n \rightarrow \infty$ то $\exists N \in \N \forall n \geq N\colon \|f - f_n\|_{\infty} < \frac{\epsilon}{2}$. Но, тогда, у $f_n$ существует неограниченная последовательность точек в которых её значение хотя бы $\frac{\epsilon}{2}$ и она не может сходиться к 0 в $+\infty$. Значит $f(t) \rightarrow 0, t \rightarrow +\infty$. Аналогично и в $-\infty$ и тогда $f \in C_{\pm\infty}(\R)$.
        \item Покажем, что $C_0(\R)$ плотно в $C_{\pm\infty}(\R)$. Пусть $f \in C_{\pm\infty}(\R)$ и $\epsilon > 0$. Поскольку $\lim\limits_{t \rightarrow \infty} f(t) = 0$, то $\exists \delta > 0 \colon \forall t\colon |t| > \delta \hookrightarrow |f(t)| < \frac{\epsilon}{2}$. Рассмотрим функцию $\tilde{f}$, которая определяется как
        \[
            \tilde{f}(t) =
            \begin{cases}
                f(t), & |t| < 2\delta \\
                f(2\delta) - \frac{f(2\delta)}{\delta}(t - 2\delta), & 2\delta \leq t \leq 3\delta \\
                f(-2\delta) + \frac{f(-2\delta)}{\delta}(t + 2\delta), & -3\delta \leq t \leq -2\delta \\
                0, & |t| \geq 3\delta.
            \end{cases}
        \]
        Функция $\tilde{f}$ -- непрерывна и имеет ограниченный носитель, а следовательно она лежит в $C_{0}(\R)$. Оценим норму разности с $f\colon$
        \[
            \|f - \tilde{f}\|_{\infty} \leq \sup\limits_{|t| \geq 2\delta} |f(t) - \tilde{f}(t)| \leq \epsilon,
        \]
        поскольку $\tilde{f}(t)$ -- убывает к нулю и $f(2\delta)$ вместе с $f(-2\delta)$ по модулю не больше $\frac{\epsilon}{2}$ и в силу нулевого предела $f$ в бесконечности при $t > \delta$ функция $f$ по модулю также не превосходит $\frac{\epsilon}{2}$. Таким образом плотность показана.
    \end{enumerate}
    Поскольку $C_0(\R)$ изометрически вкладывается в $C_{\pm\infty}(\R)$, плотно в нём и $C_{\pm\infty}(\R)$ -- полно, то $C_{\pm\infty}(\R)$ является пополнением $C_0(\R)$ с точностью до изометрического изоморфизма.
\end{solve}
\begin{task}[\S 2.7]
    Существует ли числовая функция, непрерывная в рациональных точках и разрывная в иррациональных точках отрезка $[0, 1]$?
\end{task}
\begin{solve}
    Покажем, что множество точек непрерывности произвольной функции $f$ на отрезке $[0, 1]$ является $G_\delta$-множеством. Пусть $w[f](x) = \limsup\limits_{r \rightarrow 0}\{|f(y) - f(z)| \mid y, z \in U_r(x)\}$. Рассмотрим множества $U_n = \{ x \in [0, 1] \mid w[f](x) < \frac{1}{n}\}$. Заметим, что $U_n$ -- открыты в силу полунепрерывности функции колебаний в точке сверху. В тоже время, множество точек непрерывности $f$ в точности совпадает с $\bigcap_{n = 1}^{\infty} U_n$, поскольку $f$ непрерывна в точке $x$ $\Longleftrightarrow$ $w[f](x) = 0$, что и завершает доказательство утверждения. \\
    Предположим что $f$ имеет точки непрерывности точно в $\Q \cap [0, 1]$. Тогда $[0, 1] \setminus \Q = \bigcup_{n = 1}^{\infty} K_n$ является $F_\sigma$-множеством. Заметим, что при произвольном $n \in \N$ множество $K_n$ -- нигде не плотно. Пусть это не так. Тогда, в силу замкнутости $K_n$, существует некоторый интервал, вложенный в $K_n$. Всякий непустой интервал содержит рациональную точку -- противоречие, так как тогда эта точка является точкой разрывности функции $f$. Следовательно $K_n$ -- нигде не плотно. Выберем некоторую нумерацию рациональных чисел на отрезке: $\Q \cap [0, 1] = \{q_n\}$. Тогда
    \[
        [0, 1] = \bigcup_{n = 1}^{\infty} K_n \cup \{q_n\}.
    \]
    и полное метрическое пространство представлено как счётное объединение нигде не плотных множеств -- противоречие.
\end{solve}
\begin{task}[\S 4.1]
    Доказать, что нормированное пространство полно $\Longleftrightarrow$ в нём всякий абсолютно сходящийся ряд сходится.
\end{task}
\begin{solve}
    Пусть $X$ -- нормированное пространство, $\{x_k\} \subset X$ -- последовательность и $\{S_n\}$ -- последовательность частичных сумм $\{x_k\}$, то есть $S_n = \sum_{k = 1}^n x_k$.
    \begin{enumerate}
        \item Пусть $X$ -- полно и ряд $S_n$ сходится абсолютно, то есть
        \[
            \exists \lim\limits_{n \rightarrow \infty}\sum\limits_{k = 1}^n \|x_k\| < +\infty.
        \]
        Тогда последовательность частичных сумм $\{\|x_k\|\}$ фундаментальна и $\forall \epsilon > 0 \ \exists N \in \N \  \forall m, n \geq N\colon m \geq n \wedge  \sum_{k = n}^m \|x_k\| < \epsilon$. Поскольку по неравенству треугольника
        \[
            \left\|\sum\limits_{k = n}^m x_k\right\| \leq \sum\limits_{k = n}^m \|x_k\| < \epsilon
        \]
        То и $S_n$ -- фундаментальна в $X$ и, в силу полноты $X$, сходится к некоторому $S \in X$.
        \item Пусть всякий абсолютно сходящийся ряд в $X$ сходится и $\{x_k\}$ -- фундаментальная в $X$ последовательность. Из фундаментальности следует
        \[
            \forall r \in \N \ \exists N(k) \in \N \forall m, n \geq N\colon \|x_m - x_n\| \leq 2^{-r}.
        \]
        Рассмотрим подпоследовательность вида $\{x_{N(r)}\}_{r = 1}^{\infty}$ (функцию $N(r)$ будем считать без ограничения общности строго возрастающей) и построим последовательность $\{y_r\}\colon$
        \[
            y_r = \begin{cases}
                      x_{N(1)}, & r = 1 \\
                      x_{N(r)} - x_{N(r - 1)}, & r \geq 2.
            \end{cases}
        \]
        Заметим, что $\{S_t\}$ -- последовательность частичных сумм $\{y_r\}$ -- имеет вид
        \[
            S_t = \sum\limits_{r = 1}^t y_r = x_{N(t)}.
        \]
        Покажем абсолютную сходимость ряда $\sum\limits_{r = 1}^{\infty} y_r\colon$
        \[
            \forall t \in \N \colon \sum\limits_{r = 1}^t \|y_r\| \leq \|x_{N(1)}\| + \sum\limits_{r= 1}^{t - 1} 2^{-r} \leq \|x_{N(1)}\| + 1.
        \]
        Поскольку, по предположению, всякий абсолютно сходящийся ряд сходится, то существует $x^* \in X\colon$
        \[
            x_{N(r)} \rightarrow x^*, r \rightarrow \infty.
        \]
        Но тогда и $\{x_k\}$ сходится к $x^*$ (в силу фундаментальности), а значит пространство $X$ -- полно.
    \end{enumerate}
\end{solve}
\begin{task}[\S 4.7]
    Пусть $B_1$ и $B_2$ -- шары в нормированном пространстве $(X, \|\cdot\|)$ с радиусами $r_1$ и $r_2$. Доказать, что если $B_1 \subset B_2$, то $r_1 \leq r_2$.
\end{task}
\begin{solve}
    Рассмотрим точки $x_1, x_2 \in X$ и числа $r_1, r_2\colon r_1 > r_2$. Покажем, что существует $x^* \in B_1 = B_{r_1}(x_1)\colon x^* \notin B_2 = B_{r_2}(x_2)$. Заметим, что в силу того, что $x_1 + B_{r_1}(0) = B_{r_1}(x_1)$ будем считать $x_1 = 0$ без ограничения общности. Пусть $\|x_2\| = M$. Если $M = 0$, то любой вектор с нормой из $(r_2, r_1)$ будет лежать в $B_1$, но не в $B_2$. Пусть это не так и $M \neq 0$. Рассмотрим $z_1 = \frac{r_1 - \epsilon}{M}x_2$ и $z_2 = -\frac{r_1 - \epsilon}{M}x_2$. При $\epsilon \in (0, r_1)$ эти вектора лежат в $B_1$ по определению. Пусть они лежат и в $B_2$. Тогда
    \[
        \begin{cases}
            \|z_1 - x_2\| < r_2, \\
            \|z_2 - x_2\| < r_2.
        \end{cases}
        \Rightarrow
        \begin{cases}
            |\frac{r_1 - \epsilon}{M} - 1|\|x_2\| < r_2, \\
            |\frac{r_1 - \epsilon}{M} + 1|\|x_2\| < r_2.
        \end{cases}
        \Rightarrow
        \begin{cases}
            |r_1 - \epsilon - M| < r_2, \\
            |r_1 - \epsilon + M| < r_2.
        \end{cases}
    \]
    По неравенству треугольника:
    \[
        |2r_1 - 2\epsilon| = |(r_1 - \epsilon - M) + (r_1 - \epsilon + M)| \leq |r_1 - \epsilon - M| + |r_1 - \epsilon + M| < 2r_2,
    \]
    и $\forall \epsilon \in (0, r_1)\colon$
    \[
        r_1 - \epsilon < r_2.
    \]
    Тогда, если мы перейдём к пределу при $\epsilon \rightarrow 0$, то мы получим
    \[
        r_1 \leq r_2.
    \]
    Что и завершает доказательство утверждения задачи.
\end{solve}
\begin{task}[\S 4.8]
    Пусть $B_1 \supset B_2 \supset \ldots$ -- последовательность вложенных замкнутых шаров в банаховом пространстве $(X, \|\cdot\|)$. Доказать, что $\bigcap_{n = 1}^\infty B_n \neq \emptyset$.
\end{task}
\begin{solve}
    По предыдущей задаче последовательность радиусов $\{r_n\}$ -- невозрастающая.
    Если она стремится к нулю, то задача тривиальна: в качестве точки из пересечения мы можем предел последовательности центров шаров -- он существует в силу фундаментальности последовательности центров (расстояния между любыми двумя точками при $m, n \geq N$ не превосходят $2r_N$ -- а это стремится к нулю) и полноты $X$. \\
    Пусть $\{r_n\}$ не стремится к нулю. Тогда $\exists R > 0\colon R  = \inf r_n = \lim r_n$ и $r_n = R + \epsilon_n$, где $\epsilon_n$ -- невозрастающая бесконечно малая последовательность.\\
    Оценим расстояние между центрами двух соседних в последовательности вложенных шаров $B_n(x_n)$ и $B_{n + 1}(x_{n + 1})$. Покажем, что $\|x_{n} - x_{n + 1}\| \leq r_n - r_{n + 1}$. Если $x_{n} = x_{n + 1}$, то это заведомо верно, так как $r_{n + 1} \leq r_n$. Пусть $\|x_n - x_{n + 1}\| \neq 0$. Пусть $z = x_{n + 1} + \frac{r_{n + 1}}{\|x_{n } - x_{n + 1}\|}(x_{n + 1} - x_{n})$. Такая точка лежит в $B_{n + 1}\colon$
    \[
        \|z - x_{n + 1}\| = \frac{r_{n + 1}}{\|x_n - x_{n + 1}\|}\|x_n - x_{n + 1}\| = r_{n + 1} \leq r_{n + 1}.
    \]
    А следовательно она лежит и в $B_n$. Поэтому
    \[
        \|z - x_n\| \leq r_n.
    \]
    Но $\|z - x_n\| = \|x_{n + 1} - x_n\| + r_{n + 1}$, а значит искомая оценка показана. \\
    Покажем фундаментальность последовательности центров шаров. Пусть $m, n \in \N$ и $m > n$. Тогда
    \[
        \|x_m - x_n\| \leq \sum\limits_{k = n}^{m - 1} \|x_{k + 1} - x_k\| \leq \sum\limits_{k = n}^{m - 1} r_{k} - r_{k + 1} = r_n - r_m.
    \]
    Переходя к супремуму по $m$ получим следующую оценку:
    \[
        \|x_m - x_n\| \leq \sup\limits_{m \geq n} r_n - r_m \leq r_n - R = \epsilon_n.
    \]
    Эта оценка завершает доказательство фундаментальности $\{x_n\}$. В силу полноты пространства $X$ существует точка $x^* \in X\colon x_n \rightarrow x, n \rightarrow \infty$. Она лежит в пересечении всех шаров, поскольку
    \[
        \|x_n - x^*\| = \lim_{m \rightarrow \infty} \|x_n - x_m\| \leq \lim\limits_{m \rightarrow \infty} r_n - r_m = r_n - R \leq r_n.
    \]
\end{solve}
\begin{task}[T13]
    Пусть $1 \leq p < q \leq +\infty$. Доказать, что пространство $(l_p, \|\cdot\|_q)$ является неполным, представить его в виде счётного объединения нигде не плотных множеств и построить его пополнение, состоящее из числовых последовательностей.
\end{task}
\begin{solve}
    \begin{enumerate}
        \item Для начала покажем, что $(l_p, \|\cdot\|_q)$ -- корректно определённое нормированное пространство. Пусть $f \in l_p$. Рассмотрим $A = \{n \in \N \mid |f(n)| < 1\}$. По определению $l_p\colon$
        \[
            \sum\limits_{n = 1}^{\infty} |f(n)|^p < +\infty.
        \]
        Поскольку ряд сходится, то $f(n) \rightarrow 0, n \rightarrow \infty$. В частности, это означает ограниченность $f$ и, как следствие, принадлежность $l_\infty$. Пусть $q < \infty$. При некотором $N \in \N$ верно, что $\forall n \geq N \hookrightarrow |f(n)| < 1$ -- таким образом $\N \setminus A$ конечно и сходимость суммы зависит только от элементов множества $A$. Поскольку $q > p$ то при $|f(n)| < 1$ справедливо $|f(n)|^q < |f(n)|^p$. Тогда
        \[
            \sum\limits_{n \in A} |f(n)|^q < \sum\limits_{n \in A} |f(n)|^p < +\infty.
        \]
        И $f \in l_q$.
        \item Заметим, что множество финитных последовательностей $c_0$ плотно в $l_r$ при любом $r \in [1, +\infty)$ и таким образом пополнение $(l_p, \|\cdot\|_q)$ изометрически изоморфно $(l_q, \|\cdot\|_q)$ при $q < \infty$.
        \item При $q = \infty$ пополнение по $\sup$-норме изометрически изоморфно изоморфно $c_{\pm\infty}$ -- пространству последовательностей с нулевым пределом на бесконечности. Заметим, что $(l_p, \|\cdot\|_\infty) \subset (l_\infty, \|\cdot\|_\infty)$ и $l_\infty$ -- полно, то для нахождения пополнения достаточно найти замыкание $l_p$. Поскольку $c_0$ плотно в $l_p$ и его замыкание содержит $c_{\pm\infty}$ (как просто предел усечённых последовательностей по $\sup$-норме) и $c_{\pm\infty}$ само по замкнуто (если $\{x_n\} \subset c_{\pm\infty}$ и $x_n \rightarrow x, n \rightarrow \infty$, то $\|x\|_\infty \leq \|x - x_n\|_\infty + \|x_n\|_\infty \rightarrow 0, n \rightarrow \infty$ и $x \in c_{\pm\infty}$).
        \item Пусть $F_k = \{f \in l_p \mid \|f\|_p \leq k\}$. В силу определения $l_p$
        \[
            l_p = \bigcup_{k = 1}^{\infty} F_k.
        \]
        В силу теоремы Лебега о мажорируемой сходимости $F_k$ -- замкнуты. Теперь покажем, что внутренность $F_k$ по $l_q$-норме пуста. Пусть $B_\epsilon(f) \subset F_k$, где $B_\epsilon(f) = \{h \in l_p \mid \|f - h\|_q < \epsilon\}$.

        Пусть $q < \infty$.
        Рассмотрим последовательность $h$, определяемую как
        \[
            h(n) = \dfrac{1}{n^{1/p}}.
        \]
        Её $p$-норма бесконечна:
        \[
            \|h\|_p = \left(\sum\limits_{n = 1}^{\infty} \dfrac{1}{n}\right)^{1/p} = \infty.
        \]
        А $q$-норма конечна:
        \[
            \|h\|_q^q = \sum\limits_{n = 1}^{\infty} \frac{1}{n^{q/p}} < +\infty.
        \]
        Отнормируем $h$ так, чтобы $\|\tilde{h}\|_q = \frac{\epsilon}{2}$, то есть $\tilde{h} = \dfrac{\epsilon}{2\|h\|_q}h$; $\|\tilde{h}\|_p = +\infty$. Существует $N \in \N$ такой, что $\sum\limits_{n = 1}^N |\tilde{h}(n)|^p > (2k)^p$; рассмотрим срезку $\tilde{h}$ на $N$-ом элементе:
        \[
            \hat{h}(n) = \begin{cases}
                             \tilde{h}(n), & n \leq N, \\
                             0, & n > N.
            \end{cases}
        \]
        Финитная последовательность $\hat{h}$ лежит в $l_p$. Рассмотрим $g = f + \hat{h}$. Тогда $g \in B_\epsilon(f)$, поскольку $\|\hat{h}\| \leq \frac{\epsilon}{2}$. Значит $\|g\|_p \leq k$. С другой стороны, из неравенства треугольника, $\|g\|_p \geq \|\hat{h}\|_p - \|f\|_p > k$ -- противоречие и, следовательно, $F_k$ -- нигде не плотно. Значит $(l_p, \|\cdot\|_q)$ -- первой категории Бэра и не может быть полным.


        Пусть $q = \infty$. Рассмотрим $h^M$ следующего вида:
        \[
            h^M(n) = \begin{cases}
                         \dfrac{\epsilon}{2}, & n \leq M, \\
                         0, & n > M.
            \end{cases}
        \]
        Тогда $\|h^M\|_\infty = \frac{\epsilon}{2}$ и $g^M = f + h^M \in B_\epsilon(f)$, а значит $\|g^M\|_p \leq k$. В тоже время
        \[
            \|g^M\|_p \geq \|h^M\| - k = \frac{\epsilon}{2}M^{1/p} - k. \tag{*}
        \]
        Поскольку (*) при достаточно большом $M$ будет больше $k$ то мы опять приходим к противоречию и представлению $(l_q, \|\cdot\|_\infty)$ в виде счётного объединения нигде не плотных множеств, а следовательно и неполноте $(l_q, \|\cdot\|_\infty)$.
    \end{enumerate}
\end{solve}
\begin{task}[T14]
    Пусть $1 \leq p < q \leq +\infty$. Доказать, что линейное нормированное пространство $(L^q([0, 1]), \|\cdot\|_p)$ неполно и его пополнением является пространство $(L^p([0, 1]), \|\cdot\|_p)$.
\end{task}
\begin{solve}
    \begin{enumerate}
        \item Пусть $f \in L^q([0, 1])$ и $q < \infty$. Тогда по неравенству Гёльдера
        \[
            \|f^p \cdot 1\|_1 \leq \|1\|_{1/(1 - q/p)} \cdot \|f^p\|_{q/p} = \|f^p\|_{q/p}. \tag{*}
        \]
        В тоже время
        \[
            \|f^p\|_1 = \int_X |f|^p dm, \quad \|f^p\|_{q/p} = \left(\int_X |f^q|dm\right)^{p/q}. \tag{**}
        \]
        Теперь мы подставим (**) в (*) и извлечём корень $p$-й степени и получим
        \[
            \|f\|_p \leq \|f\|_q < +\infty.
        \]
        Таким образом $L^q([0, 1]) \subset L^p([0, 1])$ в случае $q < \infty$. Поскольку на пространстве конечной меры любая степень существенно ограниченной функции интегрируема, то это включение справедливо и для $q = \infty$. Таким образом поставленная задача корректна.
        \item Заметим, что ограниченные функции плотны в $L_p([0, 1])$ при произвольном $p \in [1, +\infty)$ и их можно приблизить при помощи <<срезок>> $f^n\colon$
        \[
            f^n(t) = \begin{cases}
                         f(t), & |f(t)| < n, \\
                         0, & |f(t)| \geq n,
            \end{cases}
        \]
        по $p$-норме, поскольку по теореме о мажорируемой сходимости (с мажорантой $2|f|^p$) $\lim\limits_{n \rightarrow \infty}\|f - f^n\|^p_p = \int_{[0, 1]} \lim\limits_{n \rightarrow \infty} |f - f^n|_p^p d\mathcal{L}^1 = 0$.
        Таким образом $L^q([0, 1])$ плотно в $L^p([0, 1])$ и в силу полноты $(L^p([0, 1]), \|\cdot\|_p)$ пополнение $(L^q([0, 1]), \|\cdot\|_p)$ изометрически изоморфно $L^p([0, 1])$.
        \item Заметим, что ступенчатые функции лежат в $L_r([0, 1])$ при произвольном $r$ в силу интегрируемости ограниченных на множестве конечной меры. Поскольку всякая функция $r$-й меры может быть аппроксимирована последовательностью ступенчатых по $r$-норме, то для доказательства неполноты $(L^q([0, 1]), \|\cdot\|_p)$ нам достаточно показать строгость вложения $L^q$ в $L^p$. Рассмотрим разбиение отрезка $f$ на последовательность множеств $\{A_n\}_{n = 1}^{\infty}$ с мерами соответственно $\mathcal{L}^1(A_n) = 2^{-n}$ (оно существует, поскольку мы можем взять просто полуинтервал $[0, 0.5) \cup \{1\}$, далее брать половину непокрытого полуинтервала итеративно) и рассмотрим функцию $f\colon$
        \[
            f = \sum\limits_{n = 1}^{\infty} 2^{\frac{n}{q}} \chi_{A_n}.
        \]
        Интеграл от $|f|^q\colon$
        \[
            \int_{[0, 1]} f^qd\mathcal{L}^1 = \sum\limits_{n = 1}^{\infty} \frac{2^n}{2^n} = \sum\limits_{n = 1}^{\infty}1 = +\infty.
        \]
        И $f \notin L^q([0, 1])$. В тоже время $2^{\frac{p}{q} - 1} < 1$ (так как $\frac{p}{q} < 1$) и следовательно
        \[
            \int_{[0, 1]} f^pd\mathcal{L}^1 = \sum\limits_{n = 1}^{\infty} \left(2^{\frac{p}{q} - 1}\right)^n < +\infty.
        \]
        А значит $f \in L^p$ и утверждение задачи показано.
    \end{enumerate}
\end{solve}
\begin{task}[T15]
    Рассмотрим множество $\mathfrak{M}[0, 1]$ из задачи T11.

    Рассмотрим функцию $\rho\colon \mathfrak{M}[0, 1] \times \mathfrak{M}[0, 1] \to \R$ вида
    \[
        \rho(x, y) = \int_0^1 \dfrac{|x(t) - y(t)|}{1 + |x(t) - y(t)|}dt.
    \]
    \begin{enumerate}
        \item Доказать, что $\rho$ -- метрика в $\mathfrak{M}[0, 1]$
        \item Доказать, что метрическая топология $\tau_\rho$ совпадает с топологией $\tau_\mu$, определённой в задаче T12.
        \item Исследовать пространство $(\mathfrak{M}[0, 1], \rho)$ на полноту и сепарабельность.
    \end{enumerate}
\end{task}
\begin{solve}
    \begin{enumerate}
        \item Проверим аксиомы метрики.
        \begin{itemize}
            \item $\rho(x, x) = \int_0^1 0dt = 0$. Если же $\rho(x, y) = 0$, то и так как подынтегральная функция неотрицательна она почти всюду равна нулю, то есть $|x - y|$ почти всюду равно нулю и $x = y$ в $\mathfrak{M}[0, 1]$.
            \item Симметричность очевидна.
            \item Проверим выполнение неравенства треугольника:
            \[
                \dfrac{|x - y|}{1 + |x - y|} \leq \dfrac{|x - z|}{1 + |x - z|} + \dfrac{|z - y|}{1 + |z - y|}.
            \]
            Приведём всё к общему знаменателю и домножим на него:
            \begin{multline*}
                |x - y|(1 + |x - z||z - y| + |x - z| + |z - y|) \leq \\ \leq  |x - z|(1 + |x - y||z - y| + |x - y| + |z - y|) +  \\ + |z - y|(1 + |x - y||x - z| + |x - y| + |x - z|).
            \end{multline*}
            Вычтем совпадающие члены и получим
            \[
                |x - y| \leq |x - y||y - z||z - x| + |x - z| + |z - y| + 2|x - z||z - y|.
            \]
            Что действительно так. Поскольку неравенство справедливо при всяком $t \in [0, 1]$, то оно верно и для $\rho(\cdot, \cdot)$.
        \end{itemize}
        Значит $\rho$ -- метрика.
        \item Покажем что $\tau_\rho = \tau_\mu$, то есть проверим что $\id\colon (\mathfrak{M}[0, 1], \tau_{\rho}) \to (\mathfrak{M}[0, 1], \tau_\mu)$ является гомеоморфизмом. Достаточно показать непрерывность отображения в обе стороны. Заметим, что для $\{x_n\} \subset \mathfrak{M}[0, 1]$ и $x \in \mathfrak{M}[0, 1]$ верно следующее:
        \[
            \rho(x_n, x) \rightarrow 0, n \rightarrow \infty \Longleftrightarrow x_n \xrightarrow{\mu} x, n \rightarrow \infty.
        \]

        Пусть $\rho(x_n, x) \rightarrow 0, n \rightarrow \infty$, но сходимости по мере нет. Тогда
        \[
            \exists \epsilon > 0 \  \exists\delta > 0 \ \exists \{x_{n_k}\} \colon \mu\{t \in [0, 1] \mid |x_{n_k}(t) - x(t)| \geq \epsilon\} \geq \delta.
        \]
        Обозначим $A_k = \{t \in [0, 1] \mid |x_{n_k}(t) - x(t)| \geq \epsilon\}$ и рассмотрим $B = \limsup A_k$. Так как $\mu(B) \geq \delta$ и $t \mapsto \frac{t}{1 + t}$ -- возрастающая при $t \geq 0$, то
        \[
            \rho(x_{n_k}, x) = \int_0^1 \frac{|x_{n_k} - x|}{1 + |x_{n_k} - x|}dt \geq \int_B \dfrac{|x_{n_k} - x|}{1 +|x_{n_k} - x|}dt \geq \frac{\mu(B) \epsilon}{1 + \epsilon} > 0.
        \]
        Но тогда $\rho(x_n, x) \not\rightarrow 0, n \rightarrow \infty$ -- противоречие.

        Пусть $x_n \xrightarrow{\mu} x, n \rightarrow \infty$. Поскольку $\frac{|x_n - x|}{1 + |x_n - x|} \leq 1$ и $\frac{|x_n - x|}{1 + |x_n - x|} \leq |x_n - x|$, то при $\epsilon > 0\colon$
        \[
            \rho(x_n,x) = \int_{|x_n - x| < \epsilon} \dfrac{|x_n - x|}{1 + |x_n - x|}dt + \int_{|x_n -x| \geq \epsilon} \dfrac{|x_n - x|}{1 + |x_n - x|}dt \leq \epsilon + \mu\{|x_n - x| \geq \epsilon\} \rightarrow 0, n \rightarrow \infty, \epsilon \rightarrow +0.
        \]

        Из этого очевидно следует секвенциальная непрерывность отображения $\id$ в обе стороны. Заметим, что секвенциальная топология на $\tau_\mu$ совпадает с $\tau_\mu$, а следовательно, в силу того что $\tau_\rho$ -- first countable, то отображение непрерывно и топологии совпадают.
        \item Пространство $(\mathfrak{M}[0, 1], \rho)$ -- сепарабельно. Рассмотрим следующую систему множеств: $\left\{[\frac{j}{2^m}, \frac{j + 1}{2^m})\right\}_{j = \overline{0, 2^{m - 1}}, \ m \in \N}$. Пространство простых функций на множествах из этого разбиения с рациональными коэффициентами плотно в $\mathfrak{M}[0, 1]$.

        Из фундаментальности последовательности $\{x_n\}$ по мере следует существование фундаментальной почти всюду подпоследовательности $\{x_{n_k}\}$ и, как следствие, предела $x \in \mathfrak{M}[0, 1]$ почти всюду. Тогда и вся последовательность $x_n$ сходится к $x$ по мере.
    \end{enumerate}
\end{solve}
\section{Второе задание (неполно).}
\subsection{Компактные множества в топологических пространствах.}
\begin{task}[Т1]
    Пусть $(K, \tau)$ -- компактное топологическое пространство и $S \subset K$ его замкнутое подмножество. Доказать, что $S$ -- компакт.
\end{task}
\begin{solve}
    Пусть $\{V_\alpha\}_{\alpha \in I}$ -- открытое покрытие $S$ в индуцированной топологии. По определению индуцированной топологии $\forall \alpha \in I$ существует открытое в $K$ множество $U_\alpha\colon U_\alpha \cap S = V_\alpha$. Семейство $\{K \setminus S\} \cup \{U_\alpha\}_{\alpha \in I}$ -- открытое покрытие компакта $K$ (и действительно: так как $S$ -- замкнуто, то его дополнение открыто, а объединение всех $U_\alpha$ содержит $S$ -- таким образом это действительно покрытие) а следовательно имеет конечное подпокрытие $\{W_{i}\}_{i = 1}^n$. Заметим, что если выкинуть из этого семейства дополнение $S$ и пересечь оставшиеся множества с $S$ -- то это будет семейство открытых в $S$ множеств, объединением которых является $S$ и которое является подпокрытием изначального семейства $\{V_\alpha\}_{\alpha \in I}$, из чего следует компактность $S$.
\end{solve}
\begin{task}[Т2]
    Привести пример топологического пространства $(X, \tau)$, удовлетворяющего первой аксиоме отделимости (то есть все одноточечные множества замкнуты), и множества $S \subset X$, которое является компактом и не является замкнутым в $(X, \tau)$.
\end{task}
\begin{solve}
    Рассмотрим $\N$ с топологией ко-конечных множеств (то есть открыто пустое множество и множества с конечным дополнением). Рассмотрим множество \[S = \{2n \mid n \in \N\}.\]
    Пусть $\{U_{\alpha}\}_{\alpha \in I}$ -- произвольное его открытое покрытие. Каждое открытое множество содержит всё $\N$ за исключением, быть может, конечного числа точек. Поэтому произвольное $U_0$ из этого семейства содержит всё $S$ за исключением конечного числа точек $\{c_i\}_{i = 1}^{N}$. В силу того что $\{U_\alpha\}_{\alpha \in I}$ -- покрытие, $\forall i \in \overline{1, N}$ существует $U_i$ такое, что $c_i \in U_i$. Семейство $\{U_i\}_{i = 0}^{N}$ является искомым конечным подпокрытием $S$ (некоторые $U_i$ могут совпадать). Таким образом $S$ -- компакт.

    В тоже время $S$ не является замкнутым. И действительно: в топологии ко-конечных множеств замкнуты только конечные множества и $\N$ -- а $S$ имеет бесконечную мощность и не совпадает с $\N$, а следовательно не является замкнутым.
\end{solve}
\begin{task}[T3]
    Пусть $(X_1, \tau_1)$ и $(X_2, \tau_2)$ -- топологические пространства, отображение $f\colon X_1 \to X_2$ -- непрерывно, множество $S \subset X_1$ -- компакт в $X_1$. Доказать, что множество $f(S)$ компактно в $X_2$.
\end{task}
\begin{solve}
    Пусть $\{U_\alpha\}_{\alpha \in I}$ -- открытое покрытие $f(S)$. В силу непрерывности $f$ семейство $\{f^{-1}(U_\alpha)\}_{\alpha \in I}$ состоит из открытых множеств. Также оно является покрытием для $S$, так как $S \subset f^{-1}(f(S)) \subset \bigcup_{\alpha \in I} f^{-1}(U_\alpha)$. В силу компактности $S$ существует конечное подпокрытие $\{f^{-1}(U_i)\}_{i = 1}^n$. Но, тогда, $\{U_i\}_{i =1}^n$ -- открытое подпокрытие для $f(S)$ и, следовательно, $f(S)$ компактно.
\end{solve}
\begin{task}[T4]
    Пусть $(K, \tau)$ -- компактное пространство и $f\colon K \to \R$ -- непрерывна. Показать, что $f$ ограничена на $K$ и достигает на $K$ своих максимального и минимального значений.
\end{task}
\begin{solve}
    Образ компакта при непрерывном отображении компактен, а следовательно $f(K)$ -- компакт в $\R$. Покажем что он ограничен и замкнут.

    Ограниченность очевидна: если мы рассмотрим открытое покрытие вида $\{(n, n + 2) \mid n \in \N\}$ то после выделения из него конечного подпокрытия $f(K)$ объединение интервалов будет ограничено и сверху, и снизу -- а следовательно и $f(K)$ будет ограничено.

    Замкнутость также очевидна. Пусть $x \in \R$ -- точка прикосновения $f(K)$ и не лежит в $f(K)$. Рассмотрим семейства $\mathcal{U} = \{U(x, y)\}_{y \in f(K)}$ и $\mathcal{V} = V(x, y)\}_{y \in f(K)}$, где окрестности $U(x, y)$ и $V(x,y)$ -- непересекающиеся окрестности $x$ и $y$ соответственно (они существуют в силу хаусдорфовости $\R$). Семейство $\mathcal{V}$ является открытым покрытием $K$, а значит имеет конечное подпокрытие $\{V(x, y_i)\}_{i = 1}^N$, то есть $\bigcup_{i = 1}^{N} V(x, y_i) \supset f(K)$. С другой стороны, если мы рассмотрим $U = \bigcap_{i = 1}^{N} U(x, y_i)$, то оно будет является открытой окрестностью $x$, имеющей пустое пересечение с $f(K)$ (поскольку $U(x, y_i) \cap V(x, y_i) = \emptyset$) -- противоречие. Значит всякая точка прикосновения $f(K)$ лежит в $f(K)$ и $f(K)$ является замкнутым.

    Заметим, что $\inf f(K)$ и $\sup f(K)$ -- предельные точки $f(K)$. Они лежат в $f(K)$, а следовательно $\exists m, M \in K\colon f(m) = \inf f(K), \ f(M) = \sup f(K)$.
\end{solve}
\begin{task}[T5]
    Доказать, что топологическое пространство $[0, 1]^{[0, 1]}$ является топологическим компактом и не является секвенциальным компактом.
\end{task}
\begin{solve}
    Заметим, что отрезок $[0, 1]$ -- компакт. Следовательно, по теореме Тихонова, $\prod\limits_{t \in [0, 1]} [0, 1]_t$ -- компактно.

    В силу соображений мощности существует некоторая биекция $\phi\colon [0, 1] \to 2^{\N}$. Определим $x_n$ как
    \[
        x_n(t) = \begin{cases}
                     1, & n \in \phi(t), \\
                     0, & n \notin \phi(t).
        \end{cases}
    \]
    Пусть в $\{x_n\}$ существует подпоследовательность $\{x_{n_k}\}$, поточечно сходящаяся к некоторой функции $f$. Поскольку множество значений $x_n(t)$ при всяком $t \in [0, 1]$ состоит из двух точек -- $\{0, 1\}$, то, следовательно, последовательность $x_{n_k}(t)$ становится стационарной при достаточно большом $K(t) \in \N$.

    Рассмотрим множество индексов $I_1 = \{n_k \mid k \in \N, \ 2 \mid k\}$. Оно является бесконечным. Поскольку $I_1 \subset \N$, то существует $t_0 = \phi^{-1}(I_1)$. По определению $x_{n_k}$
    \[
        x_{n_k}(t_0) = \begin{cases}
                           1, & n_k \in I_1, \\
                           0, & n_k \notin I_1.
        \end{cases}
    \]
    Но $n_k \in I_1 \Longleftrightarrow 2 \mid k$, то есть
    \[
        x_{n_k}(t_0) = \begin{cases}
                           1, & 2 \mid k, \\
                           0, & 2 \nmid k.
        \end{cases}
    \]
    Следовательно $x_{n_k}(t_0)$ не становится стационарной ни при каком $K \in \N$, а значит у $\{x_n\}$ нет никакой сходящейся подпоследовательности и, тогда, $[0, 1]^{[0, 1]}$ не является секвенциальным компактом.
\end{solve}
\begin{task}[T6]
    Привести пример множества из топологического пространства $[0, 1]^{[0, 1]}$, которое является секвенциальным компактом и не является топологическим компактом.
\end{task}
\begin{solve}
    Рассмотрим множество $X$, определяемое как
    \[
        X = \{\Ind_{A} \mid A \subset [0, 1], \ |A| \leq \aleph_0\}.
    \]

    Покажем, что оно является секвенциальным компактом. Пусть $\{x_n\} \subset X$ -- некоторая последовательность точек из $X$. По определению $\forall n \in \N \ x_n = \Ind_{A_n}$, где $A_n$ -- некоторое не более чем счётное множество. Рассмотрим $A = \bigcup_{n = 1}^{\infty} A_n$. Оно тоже является не более чем счётным множеством, как счётное объединение не более чем счётных множеств. Выберем некоторую нумерацию $A\colon A = \{t_m \mid m \in \N\}$. При всяком $m \in \N$ множество $\{x_n(t_m)\}$ является секвенциально-компактным: оно ограничено и замкнуто в $\R$ и, следовательно, компактно -- а, поскольку, согласно критерию Фреше, для метризуемых пространств эти понятия эквивалентны, то оно является и секвенциально компактным, значит имеет некоторую сходящуюся подпоследовательность $x_{n_{k}}(t_m)$. Пусть $\{x_{n, k}\}_{n, k \in \N}$ -- последовательность сходящихся подпоследовательностей, то есть при $x_{n, k}(t_k) \rightarrow x^*(k), n \rightarrow \infty$. Рассмотрим последовательность $\{x_{n, n}\}_{n = 1}^{\infty}$. Она и будет искомой сходящейся подпоследовательностью, и будет иметь ненулевое значение не более чем на $A$ -- а следовательно будет лежать в $X$.

    Покажем, что оно не является компактом. Пусть $U(t) = \pi_t^{-1}([0, \frac{1}{2})$. Семейство открытых множеств $\{U(t)\}_{t \in [0, 1]}$ является открытым покрытием множества $X$. Пускай из него можно выделить конечное подпокрытие $\{U(t_k)\}_{k = 1}^{N}$. Рассмотрим функцию $f = \sum_{k = 1}^{N} \Ind_{\{t_k\}}$. Она очевидно лежит в $X$ -- поскольку отлична от нуля лишь в конечном числе точек. С другой стороны, она не лежит ни в каком $U(t_k)$, а следовательно не лежит и в их объединении -- следовательно, это не может быть конечным подпокрытием.
\end{solve}
\begin{task}[T7]
    Привести пример счётно-компактного топологического пространства, которое не является ни секвенциальным, ни топологическим компактом.
\end{task}
\begin{solve}
    Рассмотрим копроизведение $[0, 1]^{[0, 1]} \sqcup [0, 1]^{[0, 1]}$ и вложенное в него пространство $X \sqcup Y$, где $X = [0, 1]^{[0, 1]}$, а $Y = \{\Ind_{A} \mid A \subset [0, 1], \ |A| \leq \aleph_0\}$.

    Известно, что $X$ -- компактно, но не секвенциально компактно, а $Y$ -- секвенциально компактно, но не компактно. Поскольку как из нормальной, так и из секвенциальной компактности следует счётная компактность, то $X \sqcup Y$ -- счётно компактно.

    С другой стороны, $X \sqcup Y$ -- некомпактно: если бы оно было компактно, то, поскольку в нём с индуцированной топологии $Y$ -- замкнуто, то $Y$ было бы и компактно -- чего нет.

    Аналогично, $X \sqcup Y$ -- не секвенциально компактно. В $X$ существует последовательность, не имеющая сходящихся подпоследовательностей. Эта же последовательность точек не будет иметь сходящуюся подпоследовательность и в $[0, 1]^{[0, 1]} \sqcup [0, 1]^{[0, 1]}$.
\end{solve}
\begin{task}[T8]
    Исследовать секвенциальную и топологическую компактность топологического замыкания множества $S$ из задачи T7.
\end{task}
\begin{solve}
    Топологическим замыканием множества $S$ из задачи T7 является множество $K\colon$
    \[
        K = \{x \in \R^{(0, 1)} \mid \forall t \in (0, 1) \hookrightarrow t^2 \leq x(t) \leq t\}.
    \]

    Множество $K$ является компактным, так как $K = \prod_{t \in (0, 1)} [t^2, t]$ -- произведение компактных пространств.

    Множество $K$ не является секвенциально-компактным. Заметим, что $\forall t \in (0, 1)$ существует гомеоморфизм $\phi_t\colon [t^2, t] \to [0, 1]$. Таким образом, исходное пространство $K$ гомеоморфно $[0, 1]^{(0, 1)}$ (поскольку в силу универсального свойства произведения произвольное произведение гомеоморфизмов также будет гомеоморфизмом). При этом, поскольку из непрерывности следует секвенциальная непрерывность, то секвенциальная компактность $K$ эквивалентна секвенциальной компактности $[0, 1]^{(0, 1)}$ -- для доказательства отсутствия секвенциальной компактности которого подходит тот же пример, что и для тихоновского куба.
\end{solve}
\begin{task}[T9]
    Исследовать секвенциальную и топологическую компактность топологического замыкания множества $S$ из задачи T8.
\end{task}
\begin{solve}
    Топологическим замыканием множества $S$ из задачи T8 является множество $K\colon$
    \[
        K = \{x \in \R^{\N} \mid  \forall n \in \N \hookrightarrow n \leq x(n) \leq 2^n\}.
    \]

    Заметим, что $K = \prod_{n \in \N} [n, 2^n]$ -- и, следовательно, является компактом в топологии произведения. Поскольку $\R^{\N}$ -- метризуемо, то секвенциальная компактность в нём эквивалентна компактности, и тогда $K$ является, также, секвенциальным компактом.
\end{solve}
\subsection{Компактные и вполне ограниченные множества в метрических пространствах.}
\begin{task}[\S 3.4]
    Пусть $(X, \rho)$ -- метрическое пространство, обладающее тем свойством, что любая непрерывная на нём функция ограничена. Показать что $X$ -- компакт.
\end{task}
\begin{solve}
    Пусть $X$ не является компактом. Тогда или $X$ -- не вполне ограничено, или $X$ -- неполно.

    Пусть $X$ не является вполне ограниченным. Тогда существует $\epsilon$, для которого всякая $\epsilon$-сеть имеет бесконечную мощность. Следовательно существует последовательность точек $\{x_n\}$ такая, что $\rho(x_n, x_m) \geq \epsilon$. Рассмотрим последовательность замкнутых шаров $\overline{B}_n = \overline{B}_{\epsilon/4}(x_n)$. Определим $\phi_n$ как
    \[
        \phi_n(x) =
        \begin{cases}
            1 - \frac{4\rho(x, x_n)}{\epsilon}, & x \in \overline{B}_n \\
            0, & x \notin \overline{B_n}.
        \end{cases}
    \]
    Всякая такая функция -- непрерывна, $\phi_n(x_n) = 1$ и имеет носитель в шаре $\overline{B}_n$. Заметим, что функция
    \[
        f = \sum_{n = 1}^{\infty} n\phi_n
    \]
    является неограниченной, но при этом -- непрерывной функцией. Противоречие. Следовательно? $X$ -- вполне ограничено.

    Пусть $X$ -- неполно. Пусть $Y$ -- метрическое пополнение. Существует $x_0 \in Y \setminus X$. Рассмотрим функцию $f$, определяемую как
    \[
        f(x) = \frac{1}{\rho(x, x_0)}.
    \]
    Заметим, что в силу непрерывности метрики и отображения $t \mapsto \frac{1}{t}$ на $(0, +\infty)$, а также того, что $\rho(x, x_0) > 0$ при $x \in X$ (так как $x_0 \notin X$), то $f$ -- непрерывна. По определению метрического пополнения существует последовательность точек $X$, сходящаяся к $x_0$ в $Y$, следовательно $\rho(x_n, x_0) \rightarrow 0, n \rightarrow \infty$. Но тогда $f$ -- непрерывна, но не ограничена. Противоречие.

    Таким образом $X$ -- полно и вполне ограничено и, как следствие, является компактным пространством.
\end{solve}
\begin{task}[\S 3.8]
    Доказать что компактное метрическое пространство $(K, \rho)$ сепарабельно.
\end{task}
\begin{solve} % переделать?
    Рассмотрим семейство $\mathcal{U}_n = \left\{B_{\frac{1}{n}}(x) \mid x \in K\right\}$. Всякое такое семейство является открытым покрытием, и в силу компактности $K$ у него существует конечное подпокрытие $\tilde{\mathcal{U}}_n = \{B_{\frac{1}{n}}(x_{n, i}) \mid i = \overline{1, N(n)}\}$. Рассмотрим множество $F$ следующего вида:
    \[
        F = \bigcup_{n = 1}^{\infty} \{x_{n, i} \mid i = \overline{1, N(n)}\},
    \]
    то есть множество центров шаров, которые образуют конечные подпокрытия при всяком $n \in \N$. Оно, очевидно, счётно. Покажем его плотность в $K$.
    Пусть $x \in K$ и $\epsilon > 0$. Существует $n \in \N$ такое, что $\frac{1}{n} < \epsilon$. Поскольку $\{B_{\frac{1}{n}}(x_{n, i}) \mid i = \overline{1, N(n)}\}$ -- покрытие $K$, то $\exists j \in \overline{1, N(n)}$ такое, что $\rho(x, x_{n,j}) < \frac{1}{n}$. Следовательно $x_{n_j} \in B_\epsilon(x)$ и для всякой окрестности $x$ пересечение с $F$ непусто. Таким образом $x$ -- предельная точка $F$ и, в силу произвольного выбора $x$, $[F]_{\tau_\rho} = K$ и $K$ -- сепарабельно.
\end{solve}
\begin{task}[\S 3.10]
    Доказать, что метрический компакт $(K, \rho)$ нельзя изометрично отобразить на своё собственное подмножество.
\end{task}
\begin{solve} % переделать
    Пусть $F \subsetneq K$ и $f\colon K \to F$ -- изометрия, то есть $\forall x, y \in K \hookrightarrow \rho(x, y) = \rho(f(x), f(y))$.
    Пусть $x_0 \in K \setminus F$. Рассмотрим последовательность $\{x_n\}_{n \in \N}$, где $x_{n + 1} = f(x_{n}), x_1 = f(x_0)$ и множество $D = \{\rho(x_n, x_0) \mid n \in \N\}$.

    Если $\inf D = 0$, то существует подпоследовательность $x_{n_k}$, сходящаяся к $x_0$. Поскольку $F$ -- замкнуто (поскольку $F = f(K)$ -- компактно -- и, следовательно, замкнуто как компактное подмножество хаусдорфового пространства) и $\{x_{n_k}\} \subset F$, то $x_0 \in F$ -- противоречие.

    Пусть $\inf D = \epsilon > 0$. Тогда $\forall m, n \in \N\colon m \geq n \hookrightarrow \rho(x_m, x_n) = \rho(f^m(x_0), f^n(x_0)) = \rho(f^{m - n}(x_0), x_0) \geq \epsilon > 0$. Следовательно никакая подпоследовательность $\{x_n\}$ не является фундаментальной и, следовательно $K$ -- не секвенциально компактно. Но для метрических пространств компактность эквивалентна секвенциальной компактности -- противоречие.
\end{solve}
\begin{task}[\S 3.11]
    Доказать, что множество $M$ в $l_2$ компактно тогда и только тогда, когда оно замкнуто, ограничено и
    \[
        \forall \epsilon > 0 \ \exists n \in \N \ \forall x \in M \hookrightarrow \sum\limits_{k = n}^{\infty} |x(k)|^2 < \epsilon.
    \]
\end{task}
\begin{solve}
    Предположим, что множество $M$ -- компактно. Тогда, согласно критерию Фреше, оно полно (как метрическое пространство с индуцированной метрикой) и вполне ограничено. В частности, оно замкнуто (в силу полноты и вложенности в $l_2$) и ограничено (из вполне ограниченности следует ограниченность). Зафиксируем некоторое $\epsilon > 0$. Тогда существует некоторое $N \in \N$ такое, что существует $\epsilon$-сеть $\{x_n\}_{n = 1}^{N}$, то есть $\forall x \in M \ \exists i \in \{1, \ldots, N\}\colon \|x - x_i\| \leq \frac{\epsilon}{2}$. Заметим, что из сходимости ряда $\sum_{k = 1}^{\infty} |x_i(k)|^2 < +\infty$ следует стремление хвостов к нулю. Таким образом, $\exists I \in \N\colon \forall i \in \{1, \ldots, N\} \hookrightarrow \sum_{k = I}^{\infty} |x_i(k)|^2 \leq \frac{\epsilon^2}{4}$. Пусть теперь $x \in M$. Тогда его хвост можно оценить как
    \[
        \sqrt{\sum_{k = I}^{\infty} |x(k)|^2} \leq \sqrt{\sum_{k = I}^{\infty}|x_i(k) - x(k)|^2} + \sqrt{\sum_{k = I}^{\infty} |x_i(k)|^2} \leq \epsilon.
    \]
    Что даёт равномерную оценку на хвосты и завершает доказательство.

    Пусть теперь $M$ -- ограниченное, замкнутое множество с равномерным стремлением хвостов к нулю.
    Для доказательства компактности $M$ достаточно показать его вполне ограниченность.
    Пусть $P_n$ -- ортопроектор на подпространство, натянутое на первые $n$ базисных векторов.
    Тогда, заметим, что $P_n(M)$ -- ограничено (в силу ограниченности $M$ и $P_n$), а следовательно оно является вполне ограниченным в $\Lin\{e_1, \ldots, e_n\}$ при всяком $n \in \N$.
    Пусть $\epsilon > 0$.
    Тогда $\exists N \in \N\colon \forall x \in M \hookrightarrow \sum_{k = N}^{\infty} |x(k)|^2 < \frac{\epsilon^2}{9}$.
    Рассмотрим $\epsilon/3$-сеть в $\Lin\{e_1, \ldots, e_{N - 1}\}$ -- некоторые $\{x_i\}_{i = 1}^M$.
    Покажем, что это же множество является $\epsilon$-сетью уже во всё пространстве.
    Пусть $x \in M$.
    Тогда в $\Lin\{e_1, \ldots, e_{N - 1}\} \ \exists i \in \{1, \ldots, M\}\colon$
    \[
        \sum_{k = 1}^{N - 1} |x(k) - x_i(k)|^2 \leq \frac{\epsilon^2}{9}.
    \]
    То есть
    \[
        \|P_{N- 1}x - P_{N - 1}x_i\| \leq \frac{\epsilon}{3}.
    \]
    Тогда
    \[
        \|x - x_i\| = \|(x - P_{N - 1}x) - (x_i - P_{N - 1}x_i) + (P_{N - 1}x - P_{N - 1}x_i)\| \leq \|x - P_{N - 1}x\| + \|x_i - P_{N - 1}x_i\| + \|P_{N - 1}x - P_{N - 1}x_i\| \leq \epsilon.
    \]
    Что и завершает доказательство вполне ограниченности $M$.
\end{solve}
\begin{task}[\S 3.12]
    Пусть $(E, \rho)$ -- компактное метрическое пространство. Пусть $f\colon E \to E$, причём $\rho(f(x), f(y)) < \rho(x, y)$ при всех $x \neq y$. Показать что $f$ имеет неподвижную точку. Верно ли, что неподвижная точка -- единственна? Верно ли, что $f$ -- сжимающее отображение?
\end{task}
\begin{solve}
    Заметим, что $f$ -- непрерывная функция. И действительно: если $x_n \rightarrow x, n \rightarrow \infty$, то и $\rho(f(x_n), f(x)) < \rho(x_n, x) \rightarrow 0, n \rightarrow \infty$.

    Рассмотрим тогда отображение $\phi(x) = \rho(x, f(x))$. Оно непрерывно как композиция непрерывных и, следовательно, достигает своего минимума на компактном пространстве $E$. Тогда существует некоторая точка $x_0 \in E\colon \rho(x_0, f(x_0)) = \delta \geq 0$. Предположим, что $x_0 \neq f(x_0)$. Тогда $\rho(f(x_0), f^2(x_0)) < \rho(x_0, f(x_0)) = \delta$ -- противоречие. Следовательно $x_0 = f(x_0)$ и это искомая неподвижная точка.

    Неподвижная точка должна быть единственной, поскольку если для некоторых $x, y \in E\colon x \neq y \hookrightarrow f(x) = x, f(y) = y$, то
    \[
        \rho(x, y) = \rho(f(x), f(y)) < \rho(x, y),
    \]
    чего быть не может.

    Отображение может не быть сжимающим. Пусть $E = \{0\} \cup \{\frac{1}{n} \mid n \in \N\}$ со стандартной метрикой. Это действиетельно компактное пространство, поскольку $\frac{1}{n} \rightarrow 0, n \rightarrow \infty$, а множество точек последовательности вместе со своим пределом всегда является компактом. Определим $f$ как
    \[
        f(0) = 0, \quad f\left(\frac{1}{n}\right) = \frac{1}{n + 1}
    \]
    Тогда
    \[
        \frac{1}{(n + 1)(n + 2)} = \left|\frac{1}{n + 1} - \frac{1}{n + 2}\right| < \left|\frac{1}{n} - \frac{1}{n + 1}\right| = \frac{1}{n(n + 1)}.
    \]
    В тоже время отображение не является сжимающим:
    \[
        \frac{\left|f(\frac{1}{n}) - f(\frac{1}{n + 1})\right|}{\left|\frac{1}{n} - \frac{1}{n  +1}\right|} = \frac{n}{n + 2} \rightarrow 1, n \rightarrow \infty.
    \]
\end{solve}
\begin{task}[T10]
    Рассмотрим следующее множество $S\colon$
    \[
        S = \left\{f \in C^1[0, 1] \mid \|f\|_{C^1} \coloneq \|f\|_C + \|f'\|_C = 1 \right\}.
    \]
    \begin{enumerate}
        \item Исследовать $S$ в пространстве $(C[0, 1], \|\cdot\|_C)$  на вполне ограниченность и замкнутость.
        \item Исследовать замыкание $S$ в пространстве $(C^1[0, 1], \|\cdot\|_C)$ на вполне ограниченность и полноту.
        \item Исследовать $S$ в пространстве $(C^1[0, 1], \|\cdot\|_{C^1})$ на вполне ограниченность и замкнутость.
    \end{enumerate}
\end{task}
\begin{solve}
    \leavevmode
    \begin{enumerate}
        \item Заметим, что $S$ -- ограничено:
        \[
            \|f\|_C \leq \|f\|_C + \|f'\|_C \leq 1.
        \]
        И, более, того, $S$ -- равностепенно непрерывно:
        \[
            |f(x_1) - f(x_2)|_C = |f'(\xi(x_1, x_2)||x_1 - x_2|_C \leq |x_1 - x_2|_C.
        \]
        Следовательно, согласно теореме Арцела-Асколи, $S$ -- вполне ограничено.

        При этом $S$ не является замкнутым. Пусть $g(x) = \frac{2}{3}\left|x - \frac{1}{2}\right|$. Пусть $\{w_\epsilon\}_{\epsilon \rightarrow 0}$ -- соболевская аппроксимативная единица. Продлим $g$ на всю прямую константами (при $x < 0$ -- $g(0)$, при $x > 1$ -- $g(1)$). Заметим, что $w_\epsilon * g \xrightarrow{\|\cdot\|_C} g, \epsilon \rightarrow +0\colon$
        \begin{multline*}
            \|w_\epsilon * g - g\|_C = \sup\limits_{x \in \R} \left|\int_{\R} w_\epsilon(y)g(x - y)dy - \int_{\R} w_\epsilon(y)g(x)dy\right| = \\
            = \sup\limits_{x \in \R}\left|\int_{\R} w_\epsilon(y)(g(x - y) - g(x))dy\right| \leq \sup\limits_{x \in \R} \int_{\R} w_\epsilon(y)L|y|dy \leq \\ \leq  \sup_{x \in \R} L\epsilon = L\epsilon \rightarrow 0, \epsilon \rightarrow +0,
        \end{multline*}
        где $L$ -- константа Липшица $g$, $L = \frac{2}{3}$. При этом $w_\epsilon * g \in C^{\infty}(\R)$ и можно дать следующие оценки на нормы свёртки и её производной:
        \[
            \|w_\epsilon * g\|_C \leq \|g\|_C = \frac{1}{3}, \quad \|w_\epsilon' * g\|_C \leq \frac{2}{3}.
        \]
        И, при этом, $\|w_\epsilon * g\|_C \rightarrow \|g\|_C, \epsilon \rightarrow +0$ и $\|w_\epsilon' * g\|_C \rightarrow \frac{2}{3}, \epsilon \rightarrow +0$. Тогда, если мы рассмотрим последовательность функций $f_\epsilon$ вида
        \[
            f_\epsilon = \frac{w_\epsilon * g}{\|w_\epsilon * g\|_C + \|w_\epsilon' * g\|_C},
        \]
        то $\forall \epsilon > 0$ такие функции лежат в $S$, и они равномерно сходятся к $g$ при, например, $\epsilon \in (0, 1]\colon$
        \[
            \|f_\epsilon - g\|_C \leq C\|w_\epsilon * g - g\|_C + \left|\frac{1}{\|w_\epsilon * g\|_{C^1}} - 1\right|\|g\|_C.
        \]
        Следовательно $S$ не является замкнутым.
        \item Заметим, что та же самая последовательность функций, что и в предыдущем пункте лежит и в замыкании $S$, а её предел лежит в $C[0, 1]$, но не лежит в $C^1[0, 1]$ -- следовательно, эта последовательность будет фундаментальна, но не будет иметь предела в $\overline{S}$ -- следовательно $\overline{S}$ неполно.

        Поскольку $S$ -- вполне ограничено в $C[0, 1]$, то при всяком $\epsilon > 0$ существует конечная $\epsilon$-сеть из элементов $S$. Но, заметим, что эта же $\epsilon$-сеть будет $2\epsilon$-сетью для $\overline{S}$ в $(C^1[0, 1], \|\cdot\|_C)$ -- поскольку норма не изменилась. Следовательно, замыкание $S$ будет вполне ограниченным.
        \item Заметим, что $\|\cdot\|_{C^1}$ -- непрерывна, и, следовательно, $S = \|\cdot\|_{C^1}^{-1}(\{1\})$ является замкнутым множеством. В тоже время, согласно теореме Рисса, единичная сфера не является компактом в бесконечномерном нормированном пространстве -- следовательно, $S$ не может быть вполне ограниченным.
    \end{enumerate}
\end{solve}
\begin{task}[T11]
    Исследовать на ограниченность, вполне ограниченность и замкнутость в пространстве $c_0$ следующие множества:
    \begin{enumerate}
        \item $S_1 = \left\{x \in c_0 \mid \exists f \in L^1[0, 1]\colon \|f\|_1 \leq 1, \ \forall k \in \N \hookrightarrow x(k) = \int_{2^{-k}}^{2^{1 - k}} f(t)dt\right\}$.
        \item $S_2 = \left\{x \in c_0 \mid \exists f \in L^2[0, 1]\colon \|f\|_2 \leq 1, \ \forall k \in \N \hookrightarrow x(k) = \int_{2^{-k}}^{2^{1 - k}} f(t)dt\right\}$.
    \end{enumerate}
\end{task}
\begin{solve}
    \leavevmode
    \begin{enumerate}
        \item Заметим, что $\forall x \in S_1 \hookrightarrow$
        \[
            \|x\|_\infty = \sup_{k \in \N} \left|\int_{2^{-k}}^{2^{1 - k}}f(t)dt\right| \leq \sup\limits_{k \in \N} \int_{2^{-k}}^{2^{1- k}}|f(t)dt| \leq \sup_{k \in \N} \int_0^1 |f(t)|dt = \|f\|_1 \leq 1.
        \]
        И таким образом $S_1$ -- ограниченное множество.

        Очевидно, что $\forall x \in S_1 \hookrightarrow \|x\|_1 < +\infty\colon$
        \[
            \sum_{i = 1}^{\infty} |x_i| = \sum_{i = 1}^{\infty} \left|\int_{2^{-i}}^{2^{1 - i}}f(t)dt\right| \leq \sum_{i = 1}^{\infty}\int_{2^{-i}}^{2^{1 - i}}|f(t)|dt = \int_0^1 |f(t)|dt = \|f\|_1 \leq 1.
        \]
        То есть $S_1 \subset B^{l_1}_1(0)$.

        Пусть теперь $x \in l_1\colon \|x\|_1 \leq 1$. Тогда рассмотрим $f = \sum_{i = 1}^{\infty} 2^{i}x(i)\Ind_{[2^{-i}, 2^{1 - i}]}$. Заметим, что это корректно определённая функция из $L_1\colon$
        \[
            \int_0^1 |f(t)|dt = \sum_{i = 1}^{\infty} \int_{2^{-i}}^{2^{1-i}}2^{i}|x(i)|dt = \sum_{i = 1}^{\infty} 2^{i}|x(i)|2^{-i} = \|x\|_1 \leq 1.
        \]
        Более того, так как $\|f\|_1 = \|x\|_1 \leq 1$, то $x \in S_1$.
        Таким образом $S_1$ -- единичный шар (по $1$-норме) в $l_1 \subset c_0$.

        Пусть $\{x_n\} \subset S_1$ и $x_n \rightarrow x \in c_0$. Тогда, в частности, у нас есть поточечная сходимость. Для того, чтобы показать принадлежность $x$ множеству $S_1$ достаточно показать что $1$-норма $x$ не превосходит единицы. По лемме Фату:
        \[
            \sum_{i = 1}^{\infty} \lim_{n \rightarrow\infty} |x_n(i)| = \sum_{i = 1}^{\infty} \liminf\limits_{n \rightarrow \infty} |x_n(i)| \leq \liminf\limits_{n \rightarrow \infty} \sum_{i = 1}^{\infty} |x_n(i)| = \liminf\limits_{n \rightarrow \infty} \|x_n\|_1 \leq 1.
        \]
        Следовательно $x \in S_1$ и $S_1$ -- замкнуто.

        Заметим, что в частности, в $S_1$ содержатся базисные вектора $e_n$ и, следовательно, у нас нет равномерного убывания хвостов к нулю, а значит по критерию вполне ограниченности в $c_0$ множество $S_1$ не является вполне ограниченным.
        \item Аналогично предыдущему случаю и имея ввиду неравенство Гёльдера получаем
        \[
            \|x\|_\infty \leq \|f\|_1 \leq \|f\|_2 \leq 1.
        \]
        А значит $S_2$ -- также ограниченное множество.

        Согласно критерию вполне ограниченности в $c_0$ достаточно показать ещё равномерное убывание хвостов к нулю вместе с уже показанной ограниченностью:
        \[
            \forall \epsilon > 0 \ \exists N \in \N\colon \forall x \in S_1 \hookrightarrow \sup\limits_{k \geq N} |x(k)| \leq \epsilon.
        \]
        Поскольку
        \[
            x(k) = \langle f, \Ind_{[2^{-k}, 2^{1 - k}]}\rangle
        \]
        Согласно неравенству КБШ:
        \[
            |x(k)| = \left|\langle \Ind_{[2^{-k}, 2^{1-k}]}, f\right| \leq \|\Ind_{[2^{-k}, 2^{1-k}]}\|_2\|f\|_2 \leq 2^{-\frac{k}{2}}.
        \]
        Теперь равномерное убывание хвостов к нулю очевидно и, таким образом, $S_2$ -- вполне ограничено.

        Пусть $x_n \subset S_2$ и $x_n \rightarrow x \in c_0$. Тогда $\sup\limits_{i \in \N} |x_n(i) - x(i)| \rightarrow 0, n \rightarrow \infty$. Заметим, что $e_k = 2^{\frac{k}{2}}\Ind_{[2^{-k}, 2^{1-k}]}$ является ортонормированной системой в $L^2([0, 1])$, а $c_k(x) = 2^{\frac{k}{2}}x(k)$ для всякого $x \in S_2$ являются коэффициентами Фурье при разложении по этой системе. В силу неравенства Бесселя
        \[
            \sum_{k = 1}^{\infty} |c_k|^2 \leq \|f\|^2_2 \leq 1.
        \]
        А значит $\{c_k\}_{k = 1}^{\infty} \in l_2$. Применим это к нашей последовательности вместе с леммой Фату:
        \[
            \sum_{k = 1}^{\infty} |c_k(x)|^2 = \sum_{k = 1}^{\infty} \lim_{n \rightarrow \infty} |c_k(x_n)|^2 \leq \liminf_{n \rightarrow \infty} \sum_{k = 1}^{\infty} |c_k(x_n)|^2 \leq \liminf_{n \rightarrow \infty} 1 = 1
        \]
        Следовательно, в силу полноты $L^2([0, 1])$, $\sum_{k = 1}^{\infty} c_k(x)e_k$ сходится к некоторой функции $h \in L^2([0, 1])$. Тогда
        \[
            x(k) = 2^{-\frac{k}{2}}c_k(x) = 2^{-\frac{k}{2}}\langle g, e_k\rangle = 2^{-\frac{k}{2}} \langle g, 2^{\frac{k}{2}}\Ind_{[2^{-k}, 2^{1-k}]}\rangle = \langle g, \Ind_{[2^{-k}, 2^{1-k}]} \rangle,
        \]
        и при этом $\|g\|_2 \leq 1$ -- а следовательно $x \in S_2$ и, таким образом, $S_2$ -- замкнуто.
    \end{enumerate}
\end{solve}
\begin{task}[T12]
    Исследовать на ограниченность, вполне ограниченность и замкнутость в пространстве $C[0, 1]$ следующие множества:
    \begin{enumerate}
        \item $S_1 = \left\{f \in C[0 ,1] \mid \exists g \in C[0, 1], \|g\|_1 \leq 1, \ \forall x \in [0, 1] \hookrightarrow f(x) = \int_0^x g(t)dt\right\}$.
        \item $S_2 = \left\{f \in C[0 ,1] \mid \exists g \in C[0, 1], \|g\|_2 \leq 1, \ \forall x \in [0, 1] \hookrightarrow f(x) = \int_0^x g(t)dt\right\}$.
    \end{enumerate}
\end{task}
\begin{solve}
    \leavevmode
    \begin{enumerate}
        \item Заметим, что мы можем дать следующую оценку на $|f|\colon$
        \[
            |f(x)| = \left|\int_0^x g(t)dt\right| \leq \int_0^x |g(t)|dt \leq \int_0^1 |g(t)|dt \leq 1.
        \]
        Значит множество $S_1$ является ограниченным.

        Согласно критерию Арцела-Асколи для вполне ограниченности $S_1$ достаточно показать равностепенную непрерывность $S_1$ (вкупе с уже доказанной ограниченностью), то есть что
        \[
            \forall \epsilon > 0 \ \exists \delta > 0 \ \forall f \in S_1 \ \forall x_1, x_2 \in [0, 1]\colon |x_1 - x_2| < \delta \Rightarrow |f(x_1) - f(x_2)| \leq \epsilon.
        \]
        Пусть $g_n(x) = 2n(1 - nx) \cdot \Ind_{[0, \frac{1}{n}]}$. Тогда $\|g_n\|_1 = 1$, но для $f_n(x) = \int_0^x g_n(t)dt$
        \[
            \exists \epsilon = \frac{1}{2} \forall \delta > 0 \exists N \in \N\colon \frac{1}{N} < \delta \ \exists f_N \ \exists x_1 = 0, \ x_2 = \frac{1}{N}\colon \left|0 - \frac{1}{N}\right| < \delta \ \wedge \ \left|f(0) - f\left(\frac{1}{N}\right)\right| = 1 > \frac{1}{2}.
        \]
        Следовательно $S_1$ -- не равностепенно непрерывно и, следовательно, $S_1$ -- не вполне ограничено.

        Пусть $g = \Ind_{[0, \frac{1}{2}]}$. Рассмотрим свёртку $g$ с соболевской аппроксимативной единицей $w_\epsilon$. Покажем, что $f_\epsilon = \int_0^x w_\epsilon * g(t)dt \rightarrow f(x) = \int_0^x g(t)dt$. По определению:
        \begin{multline*}
            \|f_\epsilon - f\| = \sup\limits_{x \in [0, 1]} \left|\int_0^x w_\epsilon * g(t)dt - \int_0^x g(t)dt\right| \leq \\ \leq
            \sup_{x \in [0, 1]}\int_0^x|w_\epsilon * g - g|(t)dt \leq \sup\limits_{x \in [0, 1]} \int_0^1 |w_\epsilon * g - g|(t)dt = \\ = \int_0^1 |w_\epsilon * g - g|(t)dt \rightarrow 0, \epsilon \rightarrow +0.
        \end{multline*}
        Если существует такая непрерывная $h$, что $f(x) = \int_0^x h(t)dt$, то $f$ должна быть гладкой -- что не так. Поскольку $\|w_\epsilon * g\| \leq \|g\| \leq 1$, то множество $S_1$ является незамкнутым.
        \item Заметим, что
        \[
            |f(x)| = \left|\int_0^x g(t)dt\right| \leq \int_0^x |g(t)|dt \leq \|g\|_1 = \|g \cdot 1\|_1 \leq \|g\|_2 \cdot \|1\|_2 \leq 1.
        \]
        Значит и множество $S_2$ является ограниченным.

        Согласно критерию Арцела-Асколи для вполне ограниченности $S_2$ достаточно показать равномерную непрерывность $S_2$ (вкупе с уже доказанной ограниченностью). Поскольку
        \[
            |f(x_1) - f(x_2)| = \left|\int_{x_2}^{x_1} g(t)dt\right| = \|g\|_1^{[x_1, x_2]} \leq \|g\|_2^{[x_1, x_2]} \cdot \|1\|_2^{[x_1, x_2]} \leq \|g\|_2 \cdot \sqrt{|x_1 - x_2|} \leq \sqrt{|x_1 - x_2|}.
        \]
        Следовательно, $S_2$ -- вполне ограниченное множество.

        $S_2$ -- незамкнуто, и пример, аналогичный первому пункту здесь также подойдёт.
    \end{enumerate}
\end{solve}